% looks like distx does not define minkoshi
% looks like distx defiend twice.
% better bewone irwin hall in pick


\PassOptionsToPackage{table}{xcolor}
\documentclass[AMA,STIX2COL]{WileyNJDv5}

% ---------------- Packages ----------------
\usepackage{etoolbox}
\usepackage{geometry}
\usepackage{enumitem}
\usepackage{alltt}
\usepackage{array}
\usepackage{fontspec}
\usepackage{tabularx}
\usepackage{booktabs} 
\usepackage{ragged2e}
\usepackage{makecell}
\usepackage{xcolor}                          % <-- ADD THIS
\usepackage{array}
\newcolumntype{Y}{>{\RaggedRight\arraybackslash}X}
\usepackage[font=small]{caption}
\usepackage{tikz}
\usepackage{float} 
%\floatplacement{listing}{t}


\usepackage[most]{tcolorbox}
\usepackage{minted}
\usepackage{hyperref}
\hypersetup{hidelinks}

% ---------------- Fonts ----------------
\setmonofont{JetBrainsMono}[
  Path=./Fonts/JetBrainsMono/,
  Extension=.ttf,
  UprightFont=*-Regular,
  BoldFont=*-Bold,
  ItalicFont=*-Italic,
  BoldItalicFont=*-BoldItalic,
]

% ---------------- Spacing ----------------
\renewcommand{\baselinestretch}{0.9}

% ---------------- Colors ----------------
\definecolor{codebg}{gray}{0.98}
\definecolor{clsbg}{gray}{0.9}
\definecolor{Maroon}{RGB}{128,0,0}
\definecolor{darkyellow}{RGB}{218,165,32}

% ---------------- Custom commands ----------------
\newcommand{\IT}{{\sffamily EZR}}

% Class name highlight: \cls{Data}, \cls{Num}, etc.
\newcommand{\cls}[1]{%
  \setlength{\fboxsep}{1.5pt}%
  \colorbox{clsbg}{\sffamily\bfseries{#1}}%
}

% Solid disc marker: \disc{1} (black) or \disc[red]{2}
\newcommand{\disc}[2][black]{%
  \tikz[baseline=(c.base)]%
    \node[circle, fill=#1, draw=#1,
          inner sep=0pt, minimum size=1.2ex,
          text=white, font=\tiny\bfseries]
    (c) {#2};%
}

% Toggle for \rc markers (red <-> black)
\newtoggle{rcred}
\toggletrue{rcred}
\newcommand{\rctoggle}{\iftoggle{rcred}{\togglefalse{rcred}}{\toggletrue{rcred}}}
\newcommand{\rc}[1]{%
  \iftoggle{rcred}{\disc[red]{#1}}{\disc[black]{#1}}%
}

% "Hence" callout box for lesson banners
\newtcolorbox{hence}[1]{
  enhanced, breakable,
  colback=gray!5, colframe=black!70,
  fonttitle=\bfseries, coltitle=white,
  title=#1,
  left=5pt, right=5pt, top=4pt, bottom=4pt,
  boxrule=0.5pt, arc=2pt,
}

% ---------------- Minted setup ----------------
\captionsetup[listing]{skip=0pt,position=top}
\setminted{
  fontsize=\fontsize{6.5}{7.5}\selectfont,
  bgcolor=codebg, 
  frame=none,
  framerule=0.5pt,
  rulecolor=\color{gray},
  framesep=2mm,
  xleftmargin=0pt,
  numbersep=2mm,
  linenos,
  firstnumber=last,
  baselinestretch=0.95,
  breaklines=true,
  breakanywhere=true,
  autogobble=true,
  escapeinside=||,
}
\setminted[python]{style=trac}
\renewcommand{\theFancyVerbLine}{%
  \color{gray}\fontsize{4}{4}\selectfont\arabic{FancyVerbLine}}
\newminted[repl]{python}{linenos=false,firstnumber=last}

% Keep minted blocks from breaking across columns
\BeforeBeginEnvironment{minted}{\par\nopagebreak\minipage{\linewidth}}
\AfterEndEnvironment{minted}{\endminipage\par\noindent}

% ---------------- Wiley journal metadata ----------------
\articletype{Research Paper}
\received{Date Month Year}
\revised{Date Month Year}
\accepted{Date Month Year}
\journal{Software: Practice and Experience}
\volume{00}
\copyyear{2026}
\startpage{1}

\begin{document}

\title{Can AI be Easy? Lessons Learned from the EZR.py    Toolkit}

\author[1]{Tim Menzies}
\author[2]{Srinath Srinivasan}

\authormark{MENZIES \textsc{et al.}}
\titlemark{{Can AI be Easy? }

\address[1]{\orgdiv{Department of Computer Science},
            \orgname{North Carolina State University},
            \orgaddress{\state{North Carolina}, \country{USA}}}
\address[3]{\orgdiv{Department of Computer Science},
            \orgname{North Carolina State University},
            \orgaddress{\state{North Carolina}, \country{USA}}}

\corres{Tim Menzies, \email{timm@ieee.org}}

\abstract[Abstract]{
\renewcommand{\baselinestretch}{0.9}
Much recent press claims that 
developers no
longer need to  read code. We disagree, and
argue there is much to learn by humans who read code. 

As evidence, we present 400  lines
of EZR.py, a   Python toolkit (with no heavy dependencies) 
that implements for active learning,
classification, explanation of multi-objective optimization, regression-tree  
and text mining. EZR was built by repeatedly reading and refactoring
AI tools in order to simplify and unify them.
EZR teaches us that many seemingly different learning algorithms are nearly
the same once stripped back to their core. For example, in EZR, classical algorithms collapse to a few lines each,
and a state-of-the-art active learner fits in 35 lines.

Tested on the 120+ SE
optimization tasks in the MOOT repository, these tiny tools perform as well
as or better than state-of-the-art explanation tools (SHAP/LIME),
optimizers (SMAC/TPE), and text-mining tools- while running
500× faster, using orders of magnitude less labelled data, and building
trees from fewer than ten variables even when thousands are available.
 

We conclude that (a) AI's next step might not be larger systems but
smaller, unified ones—compact enough to be configured, critiqued, taught,
and maintained by their users; and (b) finding them means reading, very
closely, the code base of our current systems.

 EZR is available under an open source license, It can
 be installed via \textsf{pip install ezr}, which processes
 the example
 data at  \textsf{ 
github.com/timm/moot}.
}

\keywords{Active learning; explanation; multi-objective
optimization; Complement Na\"ive Bayes; software analytics;
minimalism; reproducibility; literate programming}

\jnlcitation{\cname{%
\author{Menzies T.},
\author{Srinath Srinivasan}},
\ctitle{Small AI: Lessons from Reading the EZR.py Toolkit}.
\cjournal{\it Software: Practice and Experience.}
\cvol{2026;00(00):1--20}.}

\maketitle
\section{Introduction}
\label{sec:intro}

Recently it has been argued that humans no longer need to read
code. Ryan Dahl (creator of Node.js) suggests that the era of
human-written code is ending~\cite{dahl2024}. Jensen Huang
(Nvidia CEO) advises against learning to code~\cite{huang2024}.
Elon Musk said traditional coding ends in 2026~\cite{musk2026}.
Replit's CEO told students to focus on how to talk to machines,
not how to program~\cite{masad2025}. While the phrasing varies,
the argument is the same: AI is the new compiler, and we need not
read what that compiler generates.

We disagree. Reading code is very useful. Important low-level
information ``on the ground'' (as it were) gets missed when
developers always use LLMs to ``fly over'' the details. This
paper is an extended argument that careful reading of small code
bases can offer important insights
for   widely held assumptions about
what AI and software engineering.

Here, we present 400 lines of \IT{}, a Python toolkit for
AI in software engineering. \IT{} grew across several years of
repeatedly reading and rewriting and refactoring (we read the
code, looking for similarities across distant parts of the
system, and deleting what those similarities made redundant).
The result was a small, very fast, dependency-light tool, in
which classification, clustering,  regression, explanation, multi-objective
optimization,
active learning, and text mining all sit on one shared substrate
based on three classes (\cls{Num}, \cls{Sym}, \cls{Data}) and one
update primitive (\textsf{add}).

\IT{} has been tested on the 120+ multi-objective tasks shown in
Table~\ref{mootdata} from the MOOT
repository~\cite{menzies2025moot}. For examples of these tasks
see \S\ref{details}. These
tasks are from recent papers in the software engineering
literature published, by many authors, in leading venues such as
ICSE, FSE, TSE, IST, EMSE, TOSEM, and ASE. These tasks cover a
wide range of SE problems including software configuration,
performance tuning, product line engineering, project health
forecasting, defect prediction, software testing, software
process and cost estimation, cross-domain generalization, and
text mining. To the best of our knowledge, MOOT is the largest
collection of real multi-objective SE optimization tasks yet
collected.


\begin{table*}
{\fontsize{7}{7.5}\selectfont
\setlength{\tabcolsep}{5pt}
%\renewcommand{\arraystretch}{1.15}
 

\caption{The 124 multi-objective tasks in the MOOT
repository~\cite{menzies2025moot} used in this paper. Tasks are expressed as tabular daa.
Each row
is a task category. \#y is the number of goal columns; \#x is
the range of independent attribute.
Data comes from recent papers from
the SE literature (ICSE, FSE, TSE,
IST, EMSE, TOSEM, and ASE\cite{Amiraliminimaldata,chen2026promisetune,senthilkumar2024can,chen2025accuracy,lustosa2024learning,lustossa2024isneak,chen2018beyond,nyagami_fc25_kaggle_2025,blastchar_telco_customer_churn_2025,kumarajarshi_life_expectancy_who_2025,jackdaoud_marketing_data_2022}). }

\label{mootdata}
\begin{center}
\begin{tabularx}{\textwidth}{@{} r l l l Y c c r l @{}} 
\textbf{\#}
  & \textbf{Category}
  & \textbf{Focus}
  & \textbf{File name(s)}
  & \textbf{Optimization challenge}
  & \textbf{\#y}
  & \textbf{\#x}
  & \textbf{Rows}
  & \textbf{Refs} \\
\midrule

\multicolumn{9}{@{}l}{%
  \cellcolor{gray!25}\textbf{Software systems, configuration, and tuning}} \\
24 & Soft.\ config  & PLE     & SS-[A-X]         & Minimize footprint/memory vs.\ maximize throughput   & 2--3 & 3--88   &      197--86k & \cite{Amiraliminimaldata} \\
12 & PromiseTune    & Perf.   & 7z, LLVM, BDBC   & Minimize execution time and energy consumption       & 1    & 6--35   &     864--166k & \cite{chen2026promisetune} \\
 1 & Cloud compute  & Tuning  & Apache\_AllMeas  & Balance server response vs.\ CPU/RAM load            & 1    & 9       &           192 & \cite{senthilkumar2024can} \\
 1 & Cloud compute  & Tuning  & SQL\_AllMeas     & Minimize latency/IO vs.\ maximize throughput         & 1    & 39      &      4{,}654  & \cite{chen2025accuracy} \\
 1 & Cloud compute  & Tuning  & X264\_AllMeas    & Optimize encoding parameters for PSNR/SSIM           & 1    & 16      &      1{,}153  & \cite{Amiraliminimaldata} \\
 2 & Testing        & Testing & test120, test600 & Maximize coverage while minimizing execution time    & 1    & 9       &      5{,}161  & \\

\addlinespace[3pt]
\multicolumn{9}{@{}l}{%
  \cellcolor{gray!25}\textbf{Project management and process modeling}} \\
35 & Proj.\ health  & Health  & Health-Closed    & Optimize PR rates and minimize developer churn       & 2--3 & 5       &     10{,}001  & \cite{lustosa2024learning} \\
 3 & Scrum          & Agile   & Scrum[1k--100k]  & Maximize velocity within sprint constraints          & 3    & 124     &      1k--100k & \cite{lustosa2024isneak} \\
 8 & Feature mod.   & Config  & FFM, FM-*        & Optimize clause/constraint ratio in large spaces     & 3    & 128--1k &     10{,}001  & \cite{Amiraliminimaldata} \\
 1 & nasa93dem      & Cost    & nasa93dem        & Minimize effort (person-months) vs.\ quality         & 3    & 26      &            93 & \cite{lustosa2024learning} \\
 1 & COC1000        & Risk    & coc1000          & Minimize risks vs.\ analyst expertise                & 5    & 20      &      1{,}001  & \cite{chen2018beyond} \\
 4 & POM3           & Process & pom3[a-d]        & Balance project idle rates vs.\ completion costs     & 3    & 9       &      501--20k & \cite{lustosa2024isneak} \\
 4 & XOMO           & Defects & xomo\_[flt,grd]  & Minimize defect density vs.\ total project effort    & 4    & 27      &     10{,}001  & \cite{chen2018beyond} \\

\addlinespace[3pt]
\multicolumn{9}{@{}l}{%
  \cellcolor{gray!25}\textbf{Interdisciplinary and behavioral benchmarks}} \\
 4 & Behavioral     & Behavior  & all\_players, etc.\ & Maximize user retention and predict cohort churn  & 1--3 & 26--55  &     82--17k   & \cite{nyagami_fc25_kaggle_2025} \\
 4 & Financial      & Finance   & BankChurners      & Minimize credit risk vs.\ pricing models            & 2--5 & 19--77  &      1k--20k  & \cite{blastchar_telco_customer_churn_2025} \\
 3 & Health         & Health    & COVID19, Life     & Maximize accuracy/life vs.\ readmission likelihood  & 1--3 & 20--64  &      2k--25k  & \cite{kumarajarshi_life_expectancy_who_2025} \\
 2 & RL             & Control   & A2C\_[Acr,Crt]    & Maximize reward signals vs.\ steps to convergence   & 3--4 & 9--11   &      224--318 & \\
 5 & Sales          & Market    & Marketing, socks  & Maximize ROI vs.\ minimize forecasting error        & 1--8 & 20--31  &       247--2k & \cite{jackdaoud_marketing_data_2022} \\
 3 & Misc.          & General   & auto93, Car       & Optimize multivariate trade-offs (e.g.\ MPG vs.\ HP) & 2--5 & 5--38   &       205--1k & \cite{Amiraliminimaldata} \\

\midrule
\textbf{124 } & \multicolumn{8}{@{}l}{\textbf{total}}  
\end{tabularx}
\end{center}
}
\end{table*}
\begin{table*}
\caption{This paper comments
on six common assumptions about AI for software
engineering.}
\label{truisms}
{\fontsize{8}{8.5}\selectfont
\begin{center}
\renewcommand{\arraystretch}{1.25}
\begin{tabularx}{\textwidth}{@{} l Y Y @{}}
\toprule
Section & \textbf{Common assumption} & \textbf{What this paper finds} \\
\midrule
\S\ref{sec:peek} &
AI infrastructure must be large and dependency-heavy:
real systems sit on top of \texttt{pandas}, \texttt{sklearn},
\texttt{numpy}, and a stack of supporting libraries~\cite{martinez2021}. &
Four classes (\cls{Num}, \cls{Sym}, \cls{Data}, \cls{Cols}) and
one update primitive (\textsf{add}) carry every algorithm in this
paper. A 14-line CSV reader replaces \texttt{pandas}. \\
\midrule
\S\ref{sec:famous} &
Classification, clustering, and seeded clustering are three
different algorithms that should be taugh in three different lectures and
implemented as three separate code bases. &
Classification and clustering
requires  30 lines of new code on top the 
\S\ref{sec:peek} substrate. \\
\midrule
\S\ref{trees} &
Classification trees and regression trees require different
splitting code, different leaf scoring, and different
implementations. &
43 lines of code can implement build both algorithms. And switching between them  a one-line change. \\
\midrule
\S\ref{sec:opt} &
 The 1990s metaheuristics
literature~\cite{gsat,walksat,grasp,ils} argues local search
(with restart, retry, adaptive schedules)
and
its variants dominate older methods
 &
Across 20 MOOT benchmarks, simulated annealing in its 1983
default configuration~\cite{kirkpatrick1983optimization} matches
or beats  the local-search variants we tested. Restart is a
patch for local search, not an upgrade over SA. \\
\midrule
\S\ref{sec:active} &
More labels and more features yield better models; model updates
necessitates  the complete  
retraining of  all models~\cite{hutter2011sequential,lindauer2022smac3}. &
\IT{}'s active learner reaches 85--95\% of the reference optimum
with fewer than 100 labels, uses fewer than 10 features even when
thousands are available.
The model update step of this algorithm runs  500$\times$ faster than state
of the art methods like
SMAC3~\cite{ganguly2026lowgodatalightse}. \\
\midrule
\S\ref{sec:textmine} &
Relevance filtering for retrieval-augmented generation needs
elaborate classifiers (e.g. support
vector machines) or 
LLM-grade embeddings~\cite{yu2018total,yu2018finding}. &
30 lines of complementary Bayes~\cite{rennie2003tackling},
running on the
\S\ref{sec:peek} substrate,
performs as well as prior studies
and need   roughly
one-third of the labels.\\
\bottomrule
\end{tabularx}
\end{center}}
\end{table*}

We argue that lightweight tools
like EZR, that can be quickly built and deployed,  can find important new insights.
For example,
the rest of this paper presents the six insights listed in Table~\ref{truisms}.
This table is six examples
 of the same pattern.   Our field tends to
solve problems by adding more data, more algorithms, more
parameters, and more complexity. EZR was built the other way:
by reading the existing code and removing what
turned out to be redundant.

The methodological precedent for this approach is much older than this paper.
In his 1995 book {\em
Empirical Methods for Artificial Intelligence},
Cohen~\cite{cohen1995empirical} argued that supposedly
sophisticated methods should always be benchmarked against
seemingly stupider ones --- the ``straw man'' approach to
scientific verification. We can attest that whenever we have
checked a supposedly sophisticated method against a simpler one,
there has been something useful in the simpler. More often than
not, a year later, we have switched to the simpler
approach~\cite{Nair17a,agrawal2019dodge,tantithamthavorn2016automated,
fu2016tuning}. So EZR is the latest case, but not the first.
 
The rest of this paper presents each row of
Table~\ref{truisms} in turn:
\begin{itemize}
\item
\S\ref{sec:peek} introduces the four-class substrate that is used, extensively, by the rest of the paper;
\item
\S\ref{sec:famous} implements Naïve Bayes, $k$-means, and
$k$-means++ on top of that substrate;
\item
\S\ref{trees} adds classification and regression trees;
\item 
\S\ref{sec:opt} implements and then compares simulated annealing against
local-search   on some MOOT benchmarks;
\item
\S\ref{sec:active} uses the substrate to build an active learner that runs
500$\times$ faster than   state of the art optimizers;
\item
 \S\ref{sec:textmine} discusses applications of EZR to text mining.
 \end{itemize}
 After that, the end of the paper comments on the external validity of
 our conclusions.   
\section{A Peek at the Code}\label{sec:peek}

What is the smallest set of primitives that supports
classification, clustering, optimization, regression, active
learning, and text mining? Common wisdom assumes the answer
involves \texttt{pandas}, \texttt{scikit-learn}, \texttt{numpy},
and a stack of supporting libraries~\cite{martinez2021}. This
section shows that many important parts
of those large packages can be implemted in 
80 lines of code, four classes, and one update
primitive:
\begin{itemize}
\item \cls{Num} and \cls{Sym} summarize columns of numbers and
symbols.
\item \cls{Data} holds rows together with the columns that
summarize them.
\item \cls{Cols} is a small factory that builds the columns
from a CSV header.
\item \textsf{add} is one polymorphic operation that updates any
of the above as new data arrives.
\end{itemize}
The rest of this paper builds five different AI tools on top of
those four classes. One interesting
aspect of EZR is that it is built without relying on
large third-party packages like
 \texttt{pandas} or  \texttt{sklean}.   Modern Python software
is often assembled from pre-made parts, with 85 to 97 percent
of code reused from external sources~\cite{martinez2021}.
Libraries like \texttt{pandas} are battle-tested but introduce
a large attack surface; a single poisoned update can compromise
tens of thousands of downstream packages without the primary
developer's knowledge~\cite{martinez2021}. To sidestep that
risk, \IT{} shuns \texttt{pandas} and instead ships with the
following CSV reader of 14 lines:\newpage
\begin{minted}{python}
def csv(f, ok=lambda txt: txt.partition("#")[0].split(",")):
  with open(f, encoding="utf-8") as file:
    for txt in file: |\label{csv}|
      row = ok(txt)
      if any(x.strip() for x in row):
        yield [thing(x) for x in row]

def thing(txt):
  txt = txt.strip()
  b = lambda s: {"true":1,"false":0}.get(
                  s.lower(), s)
  for f in [int, float, b]:
    try: return f(txt)
    except ValueError: pass
\end{minted}
This reader skips comments after \texttt{\#}, splits on commas,
and yields typed rows. \textsf{Thing} coerces each cell to int,
float, bool, or string in turn. 
Other packages, like \textsf{pandas} , offer the same functionality
 but at the cost
of over 400,000 lines of dependency. The 14 lines
above do just what \IT{} needs, and no more. This is the first
concrete instance of a pattern that repeats throughout the paper:
by carefully reading code, we can refactor it down to a tiny core.

\subsection{Columns: \cls{Num} and \cls{Sym}}

\IT{} distinguishes numeric and symbolic columns. Both know
their text name (\textsf{txt}), their position in the row
(\textsf{at}), and how many items they have seen (\textsf{n}).

\cls{Num}eric columns track a running mean (\textsf{mu}), a
second moment (\textsf{m2}), a standard deviation (\textsf{sd},
derived from \textsf{m2}), and a \textsf{heaven} value
indicating goal direction. By convention, columns whose names
end in ``+'' are to be maximized and columns whose names end in
``-'' are to be minimized. \textsf{heaven} is set to 1 or 0
respectively. Smaller distance to \textsf{heaven} is always
better.
\begin{minted}{python}
class Num:
  def __init__(i, txt="", at=0):
    i.txt, i.at, i.n = txt, at, 0
    i.mu = i.m2 = i.sd = 0
    i.heaven = txt[-1:] != "-" |\label{heaven}|
\end{minted}
\cls{Sym}bolic columns hold frequency counts in a \textsf{has}
dictionary:
\begin{minted}{python}
class Sym:
  def __init__(i, txt="", at=0):
    i.txt, i.at, i.n, i.has = txt, at, 0, {}
\end{minted}
The \textsf{mid} of a column is its mean (for \cls{Num}) or
mode (for \cls{Sym}):
\begin{minted}{python}
def mid(col): |\label{mid0}|
  return (col.mu if Num==type(col) else mode(col.has))

def mode(dct): return max(dct, key=dct.get)
\end{minted}
The \textsf{spread} is its standard deviation (for \cls{Num})
or entropy (for \cls{Sym}):
\begin{minted}{python}
def spread(col): |\label{spread0}|
  return (col.sd if Num==type(col) else entropy(col.has))

def entropy(dct):
  n = sum(dct.values())
  return -sum(v/n*log2(v/n) for v in dct.values())
\end{minted}
Numeric values can also be normalized to $[0,1]$ via a logistic
sigmoid clamped to $\pm3$ standard deviations. Missing values
(``?'') pass through:
\begin{minted}{python}
def norm(num, v): |\label{norml}|
  if v == "?": return v
  z = (v - num.mu) / (num.sd + 1e-32)
  return 1 / (1 + exp(-1.7 * max(-3, min(3, z))))
\end{minted}
\textsf{norm} makes rows from different scales comparable for
the distance calculations later in the paper.

\subsection{An example data file}\label{details}

With \cls{Num} and \cls{Sym} defined, the column-name
conventions \IT{} uses for tabular data become easy to state.
Names ending in ``X'' are ignored (e.g.\ row IDs or timestamps).
Names ending in ``!'' mark a symbolic class column. Names
ending in ``+'' or ``-'' mark goal columns to be maximized or
minimized. Everything else is an independent attribute. Numeric
columns start with an upper-case letter; symbolic columns do
not.

Table~\ref{moot} shows a typical example. The data come from a
configuration-tuning task in the MOOT
repository~\cite{menzies2025moot}, one of the benchmarks used
later in this paper. The header tells \IT{} everything it
needs. \texttt{Spout\_wait} starts with an upper-case letter, so
it becomes a \cls{Num}. \texttt{Spliters} and \texttt{Cntr} also
start uppercase, so they are \cls{Num}s as well.
\texttt{Throughput+} ends in ``+'', so \texttt{Throughput} is to
be maximized. \texttt{Latency-} ends in ``-'', so
\texttt{Latency} is to be minimized.

\begin{table}[!b]
\caption{An example MOOT data set: configuration tuning for a
distributed stream-processing system. Numeric column names
start with an upper-case letter. Goal column names end with
``+'' (maximize) or ``-'' (minimize). Everything else is an
independent attribute.}
\label{moot}
\begin{center}
\begin{minipage}{3in}
\normalsize{(a) Configuration tuning:}
{\fontsize{6pt}{7.2pt}\selectfont
\begin{verbatim}
x = independent values      | y = dependent values
----------------------------|-----------------------
Spout_wait, Spliters, Cntr. | Throughput+, Latency-
    10,        6,       17  |    23075,    158.68
     8,        6,       17  |    22887,    172.74
     9,        6,       17  |    22799,    156.83
[Skipped],  ...,      ...   |       ...,      ...
 10000,        1,       10  |    460.81,    8761.6
 10000,        1,       18  |    402.53,    8797.5
 10000,        1,        1  |    310.06,    9421
-------------------------------------------------------
\end{verbatim}}
\normalsize{(b) Hyperparameters for Software Health Prediction}
{\fontsize{6pt}{7.2pt}\selectfont
\begin{verbatim}
x = independent values          | y = dependent values
-------------------------------|----------------------
N_est, Min_leaf, ..., Min_imp. | MRE-, ACC+, PRED40+
   40,      1,    ...,    0.25 | 0.205, 0.826, 1.00
  120,      1,    ...,    4.75 | 0.191, 0.756, 1.00
   10,      1,    ...,    6.75 | 0.000, 0.790, 0.86
[Skipped],      ...,    ...    |   ...,   ...,   ...
   10,     11,    ...,    5.25 | 0.925, 0.391, 0.00
   90,     12,    ...,    5.00 | 0.911, 0.397, 0.00
  130,      9,    ...,    9.75 | 0.914, 0.395, 0.00
\end{verbatim}}
\end{minipage}
\end{center}
\end{table}

\subsection{Rows and tables: \cls{Data} and \cls{Cols}}

\cls{Data} holds rows together with summarized columns. Its
constructor accepts an iterator (e.g.\ from \textsf{csv}) or a
list of rows. The first row becomes the header and is handed
to \cls{Cols} to build the column objects. The rest are added
in turn:
\begin{minted}{python}
class Data:
  def __init__(i, src=None):
    src = iter(src or {})
    i.rows, i._centroid = [], None
    i.cols = Cols(next(src))
    adds(src, i)

def adds(src, it=None):
  it = it or Num(); [add(it,x) for x in src]
  return it
\end{minted}
Two usage patterns appear throughout:
\begin{repl}
>>> data1 = Data(csv(fileName))           # from disk
>>> data2 = Data([data1.cols.names])      # empty clone
\end{repl}
The second pattern (an empty \cls{Data} that mimics another's
column structure) is wrapped in a helper:
\begin{minted}{python}
def clone(data, rows=None): |\label{cloning}|
  return adds(rows or [], Data([data.cols.names]))
\end{minted}
\textsf{clone} appears nine times in the rest of
\texttt{ezr.py}. Each appearance spawns a fresh \cls{Data} that
shares its parent's schema.

\cls{Cols} is the factory that reads the header. Recall the
column-name conventions: names ending in ``X'' are ignored, ``!''
marks a symbolic class column, ``+'' or ``-'' mark goal columns
(\textsf{ys}), everything else is an independent attribute
(\textsf{xs}). Names starting with an upper-case letter become
\cls{Num}; everything else becomes \cls{Sym}.
\begin{minted}{python}
def Col(txt="", at=0):
  return (Num if txt[:1].isupper() else Sym)(txt, at)

class Cols:
  def __init__(i, names):
    i.names = names
    i.klass, i.xs, i.ys, i.all = None, [], [], []
    for j, txt in enumerate(names):
      i.all.append(col := Col(txt, j))
      if txt[-1] != "X":
        if txt[-1] == "!": i.klass = col |\label{klass}|
        role = (i.ys if txt[-1] in "+-!" else i.xs)
        role.append(col)
\end{minted}
After \cls{Cols} runs, every \cls{Data} carries three views of
its columns: \textsf{xs} (independent), \textsf{ys} (dependent
goals), and \textsf{klass} (the optional symbolic target). The
rest of \texttt{ezr.py} consults these views without re-parsing
names: decisions made once at header creation time propagate
everywhere.

\subsection{The \textsf{add} primitive}

To update a distribution, the \textsf{add} function dispatches
on the type of its first argument:
\begin{itemize}
\item For \cls{Data}, the cached centroid is zapped, the
columns are updated recursively, and the row is either
appended or removed.
\item For \cls{Cols}, it walks each column and recursively
calls \textsf{add} with the appropriate cell.
\item For \cls{Sym}, it adds (or subtracts) one from
\textsf{Sym.has}.
\item For \cls{Num}, it updates \textsf{mu}, \textsf{m2}, and
\textsf{sd} via Welford's algorithm~\cite{Knuth1998TAOCP2}, an
exact and incremental method for updating mean and standard
deviation.
\end{itemize}
\begin{minted}{python}
def add(it, v, w=1): |\label{add}|
  if Data is type(it):
    it._centroid = None |\label{zap}|
    add(it.cols, v, w) |\label{addrec1}|
    (it.rows.append if w>0 else it.rows.remove)(v)
  elif Cols is type(it):
    [add(col, v[col.at], w)
     for col in it.all] |\label{addrec2}|
  elif v != "?": # skip "don't know" values
    if Sym == type(it): it.has[v] = w + it.has.get(v, 0) |\label{symaddsub}|
    elif w < 0 and it.n <= 2:
      it.n = it.mu = it.m2 = it.sd = 0
    else: # Welford's algorithm |\label{welford}|
      it.n  += w
      delta  = v - it.mu
      it.mu += w * delta / it.n
      it.m2 += w * delta * (v - it.mu)
      it.sd = (sqrt(max(0,it.m2)/(it.n-1)) if it.n>1 else 0)
  return v
\end{minted}
Three details inside \textsf{add} are worth flagging.

\textit{First, the symmetry that makes data movement cheap.}
The \textsf{w} parameter handles both addition (\textsf{w=1})
and subtraction \mbox{(\textsf{w=-1})}, via:
\begin{minted}{python}
def sub(it, v): return add(it, v, w=-1)
\end{minted}
This symmetry matters because it lets us move a row from one
\cls{Data} to another in two constant-time operations: $\pm 1$
to a counter for symbols, or a single Welford update for
numbers. There is no retraining step, no sorting, no large-scale
data recollection. The active learner of \S\ref{sec:active}
rests on this property.

\textit{Second, the lazy centroid cache.} Every \cls{Data}
carries a \textsf{\_centroid} field that \textsf{add}
invalidates on any change (line~\ref{zap}). The \textsf{mids}
accessor lazily rebuilds it:
\begin{minted}{python}
def mids(data): |\label{mids}|
  data._centroid = data._centroid or [
                       mid(col) for col in data.cols.all]
  return data._centroid
\end{minted}
Without the cache, each \textsf{mids} call would walk every
column. With it, the call is one dictionary lookup until the
next \textsf{add}. This detail will become important when we
discuss active learning in \S\ref{sec:active}.

\textit{Third, \textsf{add} returns the value it stored.} That
return value enables a nested-call idiom used several times in
\IT{}, where the same value is removed from one distribution
and added to another in one expression; e.g.
\begin{repl}
add(rest.cols, sub(best.cols, best.rows.pop()))
\end{repl}

\subsection{Design notes}

A few stylistic points about how \texttt{ezr.py} is written
that explain the shape of the code above:
\begin{itemize}
\item Per Martin's clean-code rules\footnote{Rule 1: functions
should be small. Rule 2: they should be smaller than
that~\cite{martin2009clean}.}, \texttt{ezr.py}'s 44 functions
average $5\pm4$ lines. Each line is under 60 characters wide
so it can be included in documents like this one without
reformatting.
\item Following Hatton~\cite{Hatton:1998}, \IT{} avoids
inheritance because inheritance complicates functional
composition~\cite{GoF:1995} and debugging~\cite{Hatton:1998}.
\item Polymorphism is implemented by dispatching on type
inside functions, as seen in \textsf{add}, \textsf{mid},
\textsf{spread}, and (later) \textsf{pick}.
\item \IT{} prefers functions over methods. Methods scatter
related code across classes; functions let related code group
together. The \textsf{mid} and \textsf{spread} functions on
lines~\ref{mid0} and \ref{spread0} are examples of functions
that group functionality across different classes.
\end{itemize}

That is the entire substrate: four classes, one update
primitive, and a handful of supporting functions
(\textsf{mid}, \textsf{spread}, \textsf{norm}, \textsf{mids},
\textsf{clone}). From here on, every algorithm in this paper
is a way of asking these objects a different question. Note
that all this is achieved without including vast libraries
of other people's code.




\section{Three Classic Algorithms in 70 Lines}\label{sec:famous}

Common wisdom in AI for SE holds that classification,
clustering, and seeded clustering are three different
algorithms taught in three different lectures and implemented
as three separate code bases. As evidence of this, the reader
is invited to read the chapter headers of standard data mining textbooks
such as {\em Data Mining} by Witten et al.~\cite{Witten:2005}.

Standard ML libraries reflect
this view: scikit-learn ships each algorithm in its own
module with its own API. This section asks if that separation is necessary.
 This section shows that, on top of the
substrate from \S\ref{sec:peek}, they are nearly the same
algorithm asking slightly different questions. Naïve Bayes takes two
functions. $k$-means takes thirteen lines. $k$-means++ takes
eleven. Between them, they add 30 lines of new code and one new
primitive (\textsf{pick}, which earns its place by also serving
as the mutation operator in \S\ref{sec:opt}).

The deep reason these three algorithms collapse together is
simple: each one reduces to the same underlying question — given
some columns that summarize what we have seen, how does a new
row score against them? Naïve Bayes scores by likelihood;
$k$-means scores by distance; $k$-means++ scores by squared
distance weighted into a probability. The scoring rule changes;
the substrate that holds the summaries does not.

\subsection{Naïve Bayes}

The core of Naïve Bayes is one question: how much does a column
{\em like} a value? \textsf{like} answers that question. For a
\cls{Sym}, the answer is a smoothed frequency count from
\textsf{has} (the same dictionary \textsf{add} updates at
line~\ref{symaddsub}). For a \cls{Num}, the answer is a Gaussian
PDF using the running \textsf{mu} and \textsf{sd} that
\textsf{add} maintains via Welford (line~\ref{welford}):

\begin{minted}{python}
the.bayes.k = 1   # smoothing config

def like(col, v, prior):
  if type(col) == Sym:
    return (col.has.get(v,0) + the.bayes.k*prior
           ) / (col.n + the.bayes.k)
  sd = col.sd + 1e-32
  z  = 2*sd*sd
  return exp(-(v-col.mu)**2 / z) / sqrt(pi*z)
\end{minted}
Note what is not here: no fitting pass, no separate training
phase. The \cls{Num} and \cls{Sym} columns inside a \cls{Data}
are already a fitted Bayes model, because \textsf{add} kept them
fitted incrementally as rows arrived. \textsf{like} just reads
them. That is a non-trivial property: the substrate has erased
the distinction between fitting and using.

\textsf{likes} aggregates one row's worth of evidence. It
multiplies per-column likelihoods (in log space, for numerical
stability) with an $m$-estimate prior on class size:

\begin{minted}{python}
the.bayes.m = 2   # prior config

def likes(data, row, n_rows, n_klasses):
  prior = (len(data.rows) + the.bayes.m
          ) / (n_rows + the.bayes.m*n_klasses)
  ls = [like(col, v, prior)
        for col in data.cols.xs
        if (v := row[col.at]) != "?"]
  return log(prior) + sum(log(v)  for v in ls if v > 0)
\end{minted}
Two functions, complete with Laplace and $m$-estimate smoothing,
Gaussian numeric attributes, and missing-value handling. The
streaming test-then-train protocol that ordinarily accompanies
Naïve Bayes (classify each row with the model so far, then add
it to its true class's \cls{Data}) is a thin loop on top of
these two functions; we omit it here to keep the focus on the
algorithm itself.\footnote{The full \textsf{classify} loop is
available in \texttt{ezr.py}; it adds about a dozen lines for
streaming bookkeeping (incremental confusion matrix, a warm-up
threshold before scoring begins, and per-class \cls{Data}
maintenance).}

\subsection{$k$-means}\label{kluster}

The $k$-means clustering algorithm finds clusters by repeatedly
{\em guessing} the centroids, then {\em updating} them.  When
{\em guessing}, it labels each row with its nearest centroid
(initially, the centroids are just $k$ randomly selected rows).
{\em Update} moves centroids to the middle of the rows assigned
to them.

To talk about ``nearest centroid'' we need a distance measure. Rows in
\IT{} have mixed numeric and symbolic columns, so we need a
per-column distance that handles both, and a way to combine them
across columns. The per-column distance is \textsf{aha} (after
Aha's 1991 instance-based reasoning algorithm~\cite{aha91}); the
combiner is the Minkowski $p$-norm:\newpage

\begin{minted}{python}
the.p = 2 # config

def distx(data, r1, r2):
  return minkowski((aha(col, r1[col.at], r2[col.at])
                   for col in data.cols.xs), p=the.p)

def aha(col, u, v):
  if u == v == "?": return 1
  if Sym == type(col): return u != v
  u, v = norm(col, u), norm(col, v)
  u = u if u!="?" else (0 if v>0.5 else 1)
  v = v if v!="?" else (0 if u>0.5 else 1)
  return abs(u - v)

def minkowski(items, p=2): # p=2 means "Euclidean"  |\label{minkowski}|
  tot, n = 0, 1e-32
  for item in items:  tot, n = tot + item**p, n + 1
  return (tot/n) ** (1/p)
\end{minted}
The \textsf{aha} rules encode three   policies. If both
values are missing, their distance is maximal (Aha's
heuristic). For symbolic columns, distance is whether the two
values differ. For numeric columns, values are normalized to
$[0,1]$; if one value is missing, it is placed at the opposite
end of the known value, making the comparison deliberately
pessimistic.

With \textsf{distx} in hand, $k$-means itself is short. EZR uses
the \cls{Data} class to store each cluster: a ``cluster'' in
\IT{} is just a \cls{Data} populated by \textsf{add}:

\begin{minted}{python}
def kmeans(d, rs=None, k=10, n=10, cents=None):
  rs = rs or d.rows
  cents = cents or choices(rs, k=k)
  out = []
  for _ in range(n):
    out = [clone(d) for _ in cents] |\label{centroids}|
    for r in rs:
      add(out[min(range(len(cents)),
                  key=lambda j: distx(d, cents[j], r))], r)
    cents = [mids(kid)  # for "mids" see line |\ref{mids}|
             for kid in out if kid.rows]
  return out
\end{minted}
What is worth noting about \textsf{kmeans} is what it does
\emph{not} introduce. Every function it calls already existed
before \textsf{kmeans} was written. \textsf{clone}
(\S\ref{sec:peek}) makes the empty cluster shells. \textsf{add}
fills them. \textsf{distx} measures membership. \textsf{mids}
produces the new centroids. The clustering loop is a
fifteen-line gluing-together of pre-existing parts.

This is the substrate-reuse pattern at work. The cost of
adding a new algorithm is not the cost of building the
algorithm from scratch; it is only the cost of writing the
loop that calls the substrate.

\subsection{$k$-means++}

$k$-means++~\cite{arthur2007kmeanspp} is a heuristic for picking
the initial clusters given to $k$-means. The heuristic picks one
centroid at random, then picks each subsequent centroid with
probability proportional to its squared distance from the
nearest existing centroid. ``With probability proportional to''
is the part that needs new machinery: nothing in the substrate
so far samples from a weighted distribution.

To sample from a weighted distribution, we use the  
\textsf{pick} sampler. This is a very useful function: we use 
it here as well as when we explore
mutation-based optimization in \S\ref{sec:opt}.
That is  the same code that seeds $k$-means++ and can 
generate mutants for simulated annealing and local search. 

 To pick from \cls{Syms}, we use a 
roulette-wheel sampler: successively subtract from a probability mass,
stopping when we run out of values
 (lines~\ref{wheel0} to \ref{wheel1}).
\begin{minted}{python}
def pick(it, v=None): |\label{pick}|
  if Sym == type(it): return pick(it.has)
  if Num == type(it):
    mu = (v if v is not None and v != "?"  else it.mu) | \label{mupick} |
    lo = it.mu - 3*it.sd
    hi = it.mu + 3*it.sd
    new = mu + it.sd*2*(rand() + rand() + rand() - 1.5)
    return lo + (new - lo) % (hi - lo + 1e-32)
  if dict == type(it):
    n = sum(it.values()) * rand() | \label{wheel0} |
    for k, v in it.items():
      if (n := n - v) <= 0: break
    return k | \label{wheel1} |
\end{minted}
To pick from \cls{Num}s, we draw  frp,  a distribution centered on a supplied $\mu$ (defaulting to
the column mean) using the Irwin-Hall
sampler~\cite{johnson1995continuous}. 
To understand Irwin-Hall, we first note that a Gaussian's effective range is
$\pm 6\sigma$. The sum of $n$ uniformly distributed random numbers is approximately
normally distributed. So the sum of three uniforms, centred and
scaled, is a fast way to generate a Gaussian approximation : $2(\textsf{rand}() + \textsf{rand}() + \textsf{rand}() -
1.5)$. Mutation uses the first form (perturb around
a specific row's value); $k$-means++ uses the dict branch.

 
Once \textsf{pick} is available, $k$-means++ is just (a)~build
the weighted dict, (b)~pick from it, (c)~repeat:

\begin{minted}{python}
def kpp(d, rs=None, k=10, few=256): |\label{kpp}|
  rs = rs or d.rows
  out = [choice(rs)]
  while len(out) < k:
    t = sample(rs, min(few, len(rs)))
    ws = {i: min(distx(d, t[i], c)**2 for c in out)
          for i in range(len(t))}
    out.append(t[pick(ws)])
  return out
\end{minted}
The \textsf{few=256} parameter  (on line \ref{kpp}) is a pragmatic shortcut. On large
data sets, scanning every row to score each new centroid is
expensive; 256 random candidates give an interesting weighted
spread without the cost.

\subsection*{The lesson: one substrate, three algorithms}
The above 70 lines present in a unified manner three algorithms that
usually separate sections in a standard AI curriculum.
With the exception of \textsf{pick},
none of the three algorithms needed a new data structure or a
new support function. None needed a new training phase. Each was
a way of asking the substrate a different question:

\begin{itemize}
\item Naïve Bayes asks: how much does a column like a value?
\item $k$-means asks: which centroid is this row closest to?
\item $k$-means++ asks: which row should be the next centroid,
weighted by squared distance?
\end{itemize}
While the questions differ; the machinery underneath does not. Every
\cls{Num} and \cls{Sym} inside a \cls{Data} was already a fitted
summary. Every clustering call already had \textsf{clone},
\textsf{add}, and \textsf{mids} available do this code draws heavily on
pre=existing functionality.

This is the kind of structural similarity that becomes invisible
when humans no longer read the code. An LLM asked separately for
``Naïve Bayes,'' ``$k$-means,'' and ``$k$-means++'' will
generate three independent implementations, with three separate
data structures and three separate maintenance burdens. The
code emitted by such generators is not wrong;
but it does  misses the observation that the same substrate would
have served all three. And if the substrate is ignored and the above
is implemented three times, that  means
triple the code, triple the bugs,  and triple the maintenance. 


  \section{ Classification and Regression in 43 Lines}\label{trees}

 

 Our general theme is that with the right underlying
 architecture, building any number of AI tools becomes
 easy.
For example, in this section, we use the above code to build
classification 
 regression trees that predict the   values in any leaf of those tree.
 
 Classification trees and regression trees are nearly the same
but differ  in how they score  their leaves:
 \begin{itemize}
 \item Classification trees return symbolic   scores
 for each leaf;
 \item Regression trees return numeric scores such as \textsf{disty}.
 \end{itemize}
\begin{minted}{python}
def disty(data, row): |\label{disty}| 
  return minkowski(  # for "minkowski", see line |\ref{minkowski}|
          (abs(norm(col, row[col.at]) # for "norm", see line |\ref{norml}|
           - col.heaven)  # for "heaven", see line |\ref{heaven}|
           for col in data.cols.ys), p=the.p)
\end{minted}       
EZR's   trees are formed by recursively splitting the rows in order
 to decrease conclusion uncertainty. 
 For example,   for the following data about human beings,
 we would be certain about who was alive or dead if we split
 at   $\mathit{age}<80$
 (since, after that split {\em alive} would all $y$
 in the lower split and all $n$ in the upper split).
 \begin{repl}
 age:   [ 0, 0, 10, 11, 16, 21, 40, 80, 120, 200]
 alive: [ y, y,  y,  y,  y,  y,  y,  n,   n,   n]
\end{repl}
The \textsf{spread} function defined above at line~\ref{spread0} lets us
measure that uncertainty.
If a split on column {\em col} (using at a value {\em cut}) generates two sets of rows of
size $n_1,n_2$ with \textsf{spread} $s_1,s_2$, then the best split is
the one that minimizes the expected value of the uncertainty after
the split; i.e. $(n_1\, s_1 + n_2\, s_2)/(n_1 + n_2)$.

To assess a {\em cut},
we can the following information:
\begin{minted}{python}
def treeSplit(data, col, cut, rows,klass=disty,y=Num): 
  l_rows,r_rows,l_num,r_num=[],[],y(),y()
  for row in rows:
    v = row[col.at]
    go=v=="?" or (v==cut if Sym==type(col) else v<=cut) |\label{cutts}|
    (l_rows if go else r_rows).append(row)
    add(l_num if go else r_num,klass(data,row)) |\label{effects}|
  s=(l_num.n*spread(l_num)+r_num.n*spread(r_num))/len(rows)
  return s,col,cut,l_rows,r_rows
\end{minted}
Note how this code  uses  \textsf{spread} (from line \ref{spread0}). This is an abstract term
that covers standard deviation (for \cls{Num}s) and entropy (for \cls{Sym}s.

Line~\ref{cutts} needs a word of explanation.
Cuts from symbolic and numeric columns get tested differently. Rows go left when the test is true, right when false. For symbolic columns, true means {\em equal} to the cut (false means 
{\em not equal}). For numeric columns, true means {\em less that or equal} to  the cut (and false means {\em greater than}).

Note that, on  line \ref{effects}, \textsf{treeSplit}   uses \cls{Num} to  
collect the numeric  \textsf{disty} statistics on what happens in the left or right cut. That is, by default, EZR builds regression trees.
To ask EZR to build classification
trees, all that is required is  
\begin{repl}
klass = Sym
y = lambda row:row[data.cols.klass.at] # from line |\ref{klass}|
\end{repl}
One way to generate {\em cut}s is to try every unique value in the rows. Numeric columns can have many such values, so we instead slice the rows into equal-sized 
{\em bins} and take one 
{\em cut} per bin. After trying  various bin sizes,
 {\em bins=2} tends to work best.
(i.e. splitting numerics on the median).
\begin{minted}{python}
def treeCuts(col, rows, bins=2):
  vs = [row[col.at] for row in rows if row[col.at]!="?"]
  if not vs: return []
  if Sym == type(col): return set(vs)
  vs = sorted(vs)
  return [vs[bin*len(vs)//bins] for bin in range(1, bins)]
\end{minted}
Once we can generate and score the {\em cut}s, then the rest of tree learning is just
a recursive descent that tries to add \textsf{tree.left} and \textsf{tree.right} to the
current \textsf{tree} node.  In EZR, this recursion stops when
we can no longer cut rows into two splits of size 3 (see line \ref{stop}):
\begin{minted}{python}
the.leaf = 3 # config

def treeGrow(data, rows, klass=disty, y=Num): 
  tree = Tree(data, rows) |\label{newTree}|
  if len(rows) >= 2 * the.leaf: |\label{stop}|
    splits = (treeSplit(data, col, cut, rows, disty,y)
              for col in tree.d.cols.xs
              for cut in treeCuts(col, rows))
    if valid := [s for s in splits
        if min(len(s[3]),len(s[4])) >= the.leaf]:
      _, tree.col, tree.cut, left, right = min(
        valid, key=lambda x: x[0])
      tree.left  = treeGrow(data, left)
      tree.right = treeGrow(data, right)
  return tree
\end{minted}
 Once a tree is generated,
 to make a prediction, a new row has to be pushed down its appropriate
branch. 
\begin{minted}{python}
def treeLeaf(tree, row): 
  if not tree.left: return tree
  v = row[tree.col.at]
  go = v!="?" and (v<=tree.cut if Num==type(tree.col) 
                               else v==tree.cut)
  return treeLeaf(tree.left if go else tree.right,row)
\end{minted}
At the end that branch is a leaf containing a \cls{Data}
which can be used to make predictions.
To facilitate that, as a side-effect of making a new tree node on line~\ref{newTree}
we compute \textsf{ynum} (which is the distribution of 
the \textsf{disty}s seen in that node).
\begin{minted}{python}
class Tree: 
  def __init__(i, data, rows):
    i.d = clone(data, rows)
    i.ynum = adds((disty(data,row) for row in rows),Num())
    i.col,i.cut = None,0
    i.left=i.right=None
\end{minted}
 Later in this paper, when we discuss active learning, we will show
 that the trees learned in this manner are very effective.
 
  In support of the thesis of this paper
  (that AI is very simple
  when implemented with the right primitives) we note that EZR's
  combined classification/regression tree  generator
  requires just 43 lines of code.
One reason for this is that it makes extensive use of 
\cls{Data}, \cls{Num}, \cls{Sym}.
  
\section{Two Optimizers, One Code base, in 56 Lines}\label{sec:opt}

This section tests two truism:
\begin{itemize}
\item The first is structural: that
optimization is a different kind of problem from classification,
clustering, and regression, and that it needs its own algorithms
and its own code base. 
\item The second is empirical: that modern
optimizers (with restart, retry, adaptive schedules) outperform
older, simpler methods. The 1990s metaheuristics literature
(GSAT~\cite{gsat}, WalkSAT~\cite{walksat}, GRASP~\cite{grasp},
Iterated LS~\cite{ils}) argued the second truism explicitly — a
steady drumbeat of papers claiming clear advantage over older
methods like simulated annealing. 
\end{itemize}
The two subsections below take
these truisms in turn:
\begin{itemize}
\item
\S\ref{sec:opt}.1 shows that two optimizers (simulated
annealing and local search) share most of their code.
\item
\S\ref{sec:opt-results} then shows that simulated annealing in
its 1983 default configuration~\cite{kirkpatrick1983optimization}
matches or beats every local-search variant we tested across 20
MOOT benchmarks.
\end{itemize}
In the following implementation,  we will see  that  not only are our optimizers similar, but
they also share much of
with the same code used for classification and optimization. 
We unaware of any prior work that has shown such similarities
between optimization and classification and regression.

\subsection{(1+1) algorithms}
A (1+1) evolutionary algorithm {\em mutates} one current solution
to generate one mutant. If the mutant is ``better'' the algorithm
{\em accepts} that the current solution can be replaced. This new
current solution is used as the basis for future mutation.

Two examples of (1+1) algorithms are simulated annealing (SA) and
local search (LS). Both have their own mechanisms for escaping
dead-ends in the search space. 
Local search uses  a reset-retry policy~\cite{Crawford:1994}.
If no progress is being made, it jumps back to a fresh start with a  new random seed.

Simulated annealing~\cite{kirkpatrick1983optimization}, on the other hand, can {\em accept}
a slightly less optimal solution with a probability that drops as
the run progresses. Early on, when its ``temperature'' is high, 
SA jumps around up and down to better/worse solutions.
Later in its run, the algorithm
becomes much more discerning and will only move to better solutions.  
In the code below,
at line \ref{saxxx}, the $(1-h/b)$ term in the
denominator of the exponential is what does the cooling: as the
history count $h$ approaches the budget $b$, the denominator
shrinks, the exponential collapses, and SA stops accepting bad
moves.


The two algorithms also differ in what they {\em mutate}. Simulated
annealing mutates a random subset of attributes. Local search, on the other hand,
(a)~usually mutates any one attribute at a time; (b)~then at
probability $p$, freezes all attributes except one, then sweeps over the range
of that single attribute.

\begin{minted}{python}
def ls(d, oracle, restart=100, p=0.5, tries=20, budget=1000): 
  def accept(e, en, *_): return en < e
  def mutate(s):
    c = choice(d.cols.xs)
    for _ in range(tries if rand()<p else 1):
      s = s[:]
      s[c.at] = pick(c, s[c.at]) # for "pick" see line |\ref{pick}|
      yield s
  return oneplus1(d, # for "oneplus1", see |\ref{oneplus1}|
                  mutate, accept, oracle, budget, restart)
 
def sa(d, oracle, restart=0, m=0.5, budget=1000): 
  n = max(1, int(m * len(d.cols.xs)))
  def accept(e, en, h, b): return en<e or rand()<exp(
                                 (e-en)/(1-h/b+1E-32)) |\label{saxxx}|
  def mutate(s): 
     yield picks(d, s, n) # for "picks" see next function
  return oneplus1(d, # for "oneplus1", see |\ref{oneplus1}|
                  mutate, accept, oracle, budget, restart)

def picks(data, row, n=1): 
  s = row[:]
  for col in sample(data.cols.xs, min(n, len(data.cols.xs))):
    s[col.at] = pick(col,s[col.at]) # for "pick" see line |\ref{pick}|
  return s
\end{minted}
Simulated annealing and local search are often presented as separate
algorithms. But as seen here, both can be coded by calling an
underlying \textsf{oneplus1} function while specializing the
{\em mutate} and {\em accept} functions.

Before implementing our evolutionary algorithm, we first need to
define some way to rank mutants (in order to select the better
ones). EZR builds a surrogate {\em oracle} from a hold-out set,
then uses that surrogate to score new mutants. In the literature,
the standard way to build this {\em oracle} is a procedure followed
by (e.g.) Pfisterer~\cite{pfisterer2022yahpo} and Zela et
al.~\cite{zela2020surrogate} where:
\begin{itemize}
\item The mutators adjust the $x$ values;
\item The $y$ values of a new mutant are copied from their nearest
neighbor in 50 rows set aside at the start of the run;
\item Those $y$ values are summarized as {\em distance to heaven};
i.e. the distance of those values to the \textsf{heaven} values
for each goal calculated at line~\ref{heaven}\footnote{Some
consideration was given to building a predictive model (e.g., a
decision tree) for labeling   unseen examples. But the
only way to validate that  model would be via the Pfisterer
and Zela procedure, since our tabular data holds all   ground
truth. Since some form of nearest-neighbor testing is unavoidable,
EZR adopts it from the outset.}.
\end{itemize}
\begin{minted}{python}
d0 = Data(csv(file))   
shuffle(d0.rows)
known  = Data([d0.cols.names] + d0.rows[:50])
search = Data([d0.cols.names] + d0.rows[50:])
oracle = lambda r: oracleNearest(known, r)

def oracleNearest(data, row): 
  near = nearest(data, row)
  for col in data.cols.ys: row[col.at] = near[col.at]
  return disty(data, row)
  

def nearest(data, row): 
  return min(data.rows, 
             key=lambda r: distx(data, row, r))
\end{minted}
Finally, we can present \textsf{oneplus1}.   It runs the (1+1) loop
(mutate, score, accept, track the best) and nothing else. The
tactics used by \textsf{mutate} and \textsf{accept} are passed in
from \textsf{sa} and \textsf{ls}.
In the following code,
\textsf{restart=0} means ``never restart'' and \textsf{restart=100}
means ``restart if no improvements after 100 mutations'':
\begin{minted}{python}
def oneplus1(data, mutate, accept, oracle, |\label{oneplus1}|
             budget=1000, restart=0): 
  h, best, best_e = 0, None, 1E32
  s, e, imp = choice(data.rows)[:], 1E32, 0
  while h < budget:
    for sn in mutate(s):
      h += 1
      en = oracle(sn)
      if accept(e, en, h, budget):
        s, e = sn, en
      if en < best_e:
        best, best_e, imp = sn[:], en, h
        yield h, best_e, best
      if restart and h - imp > restart: |\label{restart}|
        s = choice(data.rows)[:]
        e, imp = 1E32, h
        break
\end{minted}
\begin{table*}[t]
\centering\scriptsize
\caption{Four treatments across 20 randomly selected MOOT benchmarks. Cells show
mean (sd) scores (from Equation~\ref{eq:wins}) over 20 repeats; budget = 1000 oracle calls. +r/-r denotes restart/retry enabled/disabled (respectively).}
\label{tab:opt}
\begin{tabular}{rrrrrrrrl}
wins & ls-r & ls+r & sa-r & sa+r & \#rows & \#x & \#y & file \\
\midrule
2 &  30 (75) &  89 (17) &  98 (9)  &  98 (9)  & 86058 & 11 & 2 & SS-X \\
2 &  45 (62) &  97 (9)  & 100 (0)  & 100 (0)  &  4653 & 38 & 1 & SQL\_AllMeasurements \\
2 &  50 (47) &  91 (18) & 100 (0)  & 100 (0)  &  4608 & 21 & 2 & SS-U \\
3 &  54 (84) & 100 (0)  & 100 (0)  & 100 (0)  &   864 & 17 & 3 & SS-M \\
2 &  62 (16) &  80 (11) &  92 (11) &  92 (11) & 53662 & 17 & 2 & SS-N \\
2 &  63 (43) &  96 (12) & 100 (0)  & 100 (0)  &  1023 & 11 & 2 & SS-P \\
1 &  74 (20) &  93 (12) &  88 (13) &  92 (13) &   196 &  3 & 2 & SS-G \\
1 &  76 (19) &  96 (11) &  94 (13) &  97 (8)  &   196 &  3 & 2 & SS-F \\
3 &  77 (40) & 100 (0)  & 100 (0)  & 100 (0)  &  1152 & 16 & 1 & X264\_AllMeasurements \\
3 &  81 (49) & 100 (0)  & 100 (0)  & 100 (0)  &  3840 &  6 & 2 & SS-S \\
1 &  83 (21) &  98 (9)  &  96 (11) &  95 (13) &   196 &  3 & 2 & SS-D \\
3 &  92 (25) & 100 (0)  & 100 (0)  & 100 (0)  &   196 &  3 & 1 & wc+wc-3d-c4-obj1 \\
3 &  92 (34) & 100 (0)  & 100 (0)  & 100 (0)  &  5184 & 12 & 2 & SS-T \\
3 &  93 (23) & 100 (0)  & 100 (0)  & 100 (0)  &   192 &  9 & 1 & Apache\_AllMeasurements \\
2 &  94 (13) &  99 (5)  & 100 (0)  & 100 (0)  &   259 &  4 & 2 & SS-H \\
3 &  95 (23) & 100 (0)  & 100 (0)  & 100 (0)  &   196 &  3 & 1 & wc+sol-3d-c4-obj1 \\
3 &  98 (11) & 100 (0)  & 100 (0)  & 100 (0)  &  3840 &  6 & 2 & SS-J \\
4 & 100 (0)  & 100 (0)  & 100 (0)  & 100 (0)  &   196 &  3 & 1 & wc+rs-3d-c4-obj1 \\
4 & 100 (0)  & 100 (0)  & 100 (0)  & 100 (0)  &  3840 &  6 & 1 & rs-6d-c3\_obj1 \\
4 & 100 (0)  & 100 (0)  & 100 (0)  & 100 (0)  &  2866 &  6 & 1 & sol-6d-c2-obj1 \\
\midrule
  &  78      &  97      &  98      &  99      &       &    &   & mean 
\end{tabular}
\end{table*}
Note that:
\begin{itemize}
\item The \textsf{restart} parameter jumps the search back to a
fresh random starting row if no improvement has appeared in the
last \textsf{restart} evaluations.  
\item \textsf{sa} calls \textsf{oneplus1} with \textsf{restart=0}
which, as can be seen on line~\ref{restart}, disables restarts for
that algorithm.
\end{itemize}
Note also that, as promised at the start of this section, not only
do these two optimizers share much of the same codebase, but this
code also reuses large parts of the clustering code:
\begin{itemize}
\item Our optimizers needed a surrogate oracle to assess new
mutants. We supplied this with a nearest neighbor algorithm based
on the clustering code  described in \S\ref{kluster}.
\item When mutating numerics, our optimizer uses distribution
knowledge of the standard deviation to decide what counts as a
reasonable mutation size. This information was already collected
by the \cls{Num} class, which we were using for the surrogates.
\end{itemize}



 
\subsection{Experiments with (1+1)}\label{sec:opt-results}

 

Code similarities aside, another reason to use EZR is that it is
very easy to reorganize such a small code base to conduct insightful
experiments. For example, in this section we comment
on the relative merits of restart/retry mechanisms.


Why is it interesting to explore this issue?
In the 1990s, LS with all its variants (GSAT~\cite{gsat},
WalkSAT~\cite{walksat}, GRASP~\cite{grasp}, Iterated LS~\cite{ils})
represented a dominant paradigm in metaheuristics research—a 
steady drumbeat of papers claiming clear advantage over older 
methods like SA. A reader who
read that literature would expect LS to dominate here.

However, our results with EZR on the MOOT data
(Table~\ref{mootdata}) tell a different story. This section
details a recurring observation: many long-standing truisms
in optimization dissolve when confronted with modern
empirical benchmarks.




To compare LS and SA,  we need data and some goal
within that data
For data, we selected 20 data sets at random from the MOOT data
 (recall
that examples of those tables were displayed in 
Table~\ref{moot}). 


Turning now to what goal to explore, within any MOOT table
of data there is one row with
the smallest \textsf{disty}. We call this the {\em reference
optimum} $d^*$.\footnote{According to the
empirical algorithms literature~\cite{McGeoch2012,cohen1995empirical},
the 
{\em reference optimum} is the best solution observed in the dataset. This is not necessarily the true global optimum
since in most real engineering problems,
the true optimum is unknown, so the reference optimum is the
strongest available standard}.   A   baseline $d_0$ is set to
the median \textsf{disty} of the rows in that table.
In order to compare optimization results across multiple
MOOT data sets, we can report the best row found as some
ratio of $d_0 - D^*$: 
\begin{equation}\label{eq:wins}
  \text{score} = 100 \cdot \left( 1 - \frac{\text{disty(best)} - d^*}{d_0 - d^*} \right)
\end{equation}
Here, the {\em score}$=0$  means the optimizer did no better than a random pick and {\em score}$=50$
means it closed half the gap, and 
{\em score}$=100$ means it found the reference
optimum.

(Aside: for readers familiar with reinforcement learning,
we note that Equation~\ref{eq:wins} is the
complement of {\em regret}, normalized by the achievable
improvement range.)


 For each of the 20 MOOT benchmarks used in this study, we shuffle the
rows, set aside 50 as the surrogate's known set, and run 
LS and SA with restart/retry enabled/distabled
for a budget of 1000 oracle calls. We repeated
this 20 times per (file, treatment) pair, recording the score of
the best mutant found.

 Table~\ref{tab:opt} reports mean and standard
deviation across the 20 repeats. 
In those results, two patterns are noteworthy:
\begin{itemize}
\item
\textit{Restart is a patch for LS, not an upgrade over SA.}
Removing restart from LS is catastrophic: the \textsf{ls-r} score mean
drops to 78, and on the worst rows it lands at 30, 45, 50, with
standard deviations of 75, 62, 47.   Adding restart back fixes
LS (\textsf{ls+r} mean of 97), but does not push it past SA. SA
without restart already averages 98, and adding restart to SA
moves the mean by one point. 
To say that another way,  
SA's stochastic acceptance is doing
the same job as restart.
\item
\textit{The variance is the story, not the mean.} The \textsf{ls-r}
column has standard deviations of 75, 84, 62, 49 on its hardest
rows. The other three columns sit at 0--13 sd on those same rows.
For a practitioner, this means \textsf{ls-r} is not just worse on
average, it is unpredictable in a way the others are not.  
\end{itemize}
The headline for a practitioner is as follows:
 reset/retry should be not be used uncritically.
Rather, in new domains, it would be useful to comparatively
assess the merits of reset/retry against other approaches.
If you are running local search,
you need restart, and even with it you do not pull ahead of SA.
If you are running SA, restart is a rounding error. The simplest
working choice is \textsf{sa-r}: simulated annealing with the
default \textsf{restart=0}, the same configuration the algorithm
has had since 1983.

These results come with the obvious caveat: it only holds
for the data studied here. Further experimentation, on more data,
is required.  Tools like EZR encourage such experimenation
due to their simplicity.  

In summary: (a)~optimization is not its own
separate algorithmic kingdom (since it shares so much with regression and classification), 
and (b)~1983 SA still holds its ground against
the 1990s LS variants that claimed to surpass it.

\section{Active Learning in 80 Lines }\label{sec:active}

This section tests  truisms about learning
from labeled data. 
When learning some model $y=f(x)$ for examples of $(x,y)_1, (x_y)_2,$ etc,
it is often assumed that:
\begin{itemize}
\item
The  more labels for the $y$ attributes, the better the models;
\item 
Using 
more  $x$ features yields better models;
\item
Retraining models is always  a slow process.
\end{itemize}
This section offers counter examples to all these points.
% The active learner described
% below pushes back on all three. The \textsf{rebalance} routine
% below replaces ensemble retraining with an O(1) row swap,
% running 500$\times$ faster than 
% SMAC3~\cite{hutter2011sequential,ganguly2026lowgodatalightse}.
% A parameter study across 100{,}000 trials shows fewer than 100
% labels suffice to reach 85--95\% of the reference optimum. And
% the regression trees grown from those labels use fewer than 10
% features even when thousands of features are available.

\subsection{Why Active Learning?}
The prior results with (1+1) algorithms
assumed  that  we could  obtain   $y$ labels of up to 1000 rows.  
In SE, that assumption is unwise:
\begin{itemize}
\item
{\bf Labeling by human experts}  is possible but slow and error-prone when rushed~\cite{easterby1980design}, often taking hours for just a few cases~\cite{KingtonAlison2009,lustosa2024learning,valerdi2010heuristics}.
\item
{\bf Historical logs} provide large label sets but can be unreliable; e.g., 90\% of technical debt “false positives” were incorrect~\cite{yu2020identifying}; similar issues occur in security~\cite{wu2021data}, static analysis~\cite{kang2022detecting}, and defect data~\cite{shepperd2013data}.
\item
{\bf Automated labeling} also faces limits: regex-based heuristics are crude~\cite{kamei2012large}; LLMs are only assistive (but not authoritative).
Even in domains with naturally occurring oracles (e.g. compile with certain Makefile settings; run the full test suite), these can be extremely slow. For example, an exhaustive
exploration of the 11 parameters of the \textsc{x264} video encoder took 1000+  hours to terminate~\cite{DBLP:conf/wosp/ValovPGFC17}. Such costs mean that, in practice, the number of evaluations is often limited to just a few dozen~\cite{DBLP:journals/tse/Nair0MSA20,DBLP:conf/icse/0003XC021}.
\end{itemize}
% Finding authoritative labels is a  problem
% that has plagued much recent work~\cite{ahmed2025can,h2022,tu2019better,yu2020identifying,wu2021data,kamei2012large,herbold2022problems}.   Labels can be collected from domain experts but
%  label quality degrades when those experts are rushed~\cite{easterby1980design}. Careful interviews are slow, requiring hours to label just a few examples~\cite{KingtonAlison2009, lustosa2024learning, valerdi2010heuristics}\footnote{Similar rates were seen  in manual knowledge acquisition research from the 1980s where `5-10 expert rules per day'' was the norm~\cite{menzies1992expert}.}.
%  Historical logs  offer many  labels, but many of those are error-prone (e.g. 90\% of technical debt ``false positives'' were found to be incorrect\cite{yu2020identifying}; similar errors plague security\cite{wu2021data}, static analysis~\cite{kang2022detecting}, and defect datasets~\cite{shepperd2013data}).
%  Automated labeling models can be applied, but these are  often problematic. Simplistic regex methods rely on dubious heuristics~\cite{kamei2012large}, while recent MSR'25 studies found that LLMs are only assistive, not authoritative~\cite{ahmed2025can}. Even widely-used tools like SZZ~\cite{Sliwerski05} have many known issues, with some doubting that they will ever be authoritative~\cite{herbold2022problems}. Ultimately, all automated methods require certification, making them dependent on experts or logs.
 When we cannot trust many labels, we must do what we can with very few labels. Enter {\em active learning}~\cite{settles2008analysis}. 
 
 A {\em passive learner} trains on whatever data it is given (the Naive Bayes and
 $k$-means algorithms described above are passive learners). An {\em active
learner}, on the other hand,  picks its own training set. Specifically, an active learner reflects on the
model built so far to find the  most informative    example to label next. 
It then updates its model and the process repeats. In practice,
this means the active learner labels very few rows since it is 
skipping over all the  examples that are noisy, redundant, or
uninformative.  

\subsection{Taming the Cost of Model Update}

The cost of active learning is in the model updates. New labels arrive
one at a time. After each arrival, the model has to be updated. If
the model is an ensemble of regression trees (as in
SMAC3~\cite{hutter2011sequential}), this update means rebuilding all the trees.

EZR's active learner 
was inspired by SMAC3. For every part of SMAC3, we asked if   analogous
functionality could be achieved in a simpler way.
For example,   SMAC maintains many models in its ensemble.
Why not try just two models?  We tried:
\begin{itemize}
\item A \cls{Data} called \textsf{best} holdimg the top $\sqrt{N}$
(where ``best'' is defined by \textsf{disty} from line~\ref{disty}), where $N$ is the total number of
labeled rows so far.
\item A second \cls{Data} called \textsf{rest} holds the
remaining $N - \sqrt{N}$.
\item When a new label arrives, it enters \textsf{best}. If
  \textsf{best} grows larger than  $\sqrt{N+1}$, the worst row in
\textsf{best} is demoted to \textsf{rest}.
\end{itemize}
EZR implements this demotion using the same \textsf{add} function used
above for $k$-means and our Naive Bayes classifier:
\begin{minted}{python}
def rebalance(best, rest, lab): 
  if len(best.rows) > sqrt(len(lab.rows)):
    best.rows.sort( |\label{sorting}|
     key=lambda row:disty(lab,row)) # for "disty" see line |\ref{disty}|
    add(rest.cols, 
       sub(best.cols,
         best.rows[-1]))
    rest.rows.append(best.rows.pop())
\end{minted}
The only computationally expensive part of this loop
is the call to \textsf{best.rows.sort} on line \ref{sorting}. But even this
is not particularly expensive since \textsf{best} is bounded to $\sqrt{N}$;
i.e.\ 
this sort only has to rearrange a handful of items.  As to the larger \textsf{rest} sets,  EZR never needs to resort those rows.

Ganguly reports  that this use of   \textsf{rebalance}
makes EZR run
 500 times faster than SMAC3~\cite{ganguly2026lowgodatalightse}.  
The reason for this is simple: 
 SMAC3 has to rebuild
all its trees each time a new row is labeled.
EZR, on the other hand, just jumps a row from
{\em best} to {\em rest}.
As seen above (at line \ref{add}), adding or subtracting a row for a \cls{Data} requires very little computation.  

Having tamed the cost of model update, we can now define the active learning loop that calls
\textsf{rebalance}.

 \subsection{ From Active Learning to Models}
In order to test the generality
 of a learning process,
 it is best to learn a model
 from some training data, then apply it to
 test data.
 Hence, in order to test the code presented
 in the last section,
 we have to commit
 to a specific  model format (in order to transfer knowledge from   training to testing). But in that process, what model format should be used?

% One format supported by EZR is the tin

%  Rayegan et al.~\cite{Amiraliminimaldata} recommend  
%  {\em tiny regression trees} as
%  an appropriate model format. 
%  EZR's
%  active learning methods select
%  just a few dozen most
%  informative examples. Why build complex models
%  that no one understands, argues
%  Rayegan et al., when we can just
%  generate  tiny trees from those few dozen examples? 





 One model format supported by EZR is the {\em tiny regression tree}
 such as shown in 
Figure~\ref{htree}.
That tree was grown using the code from \S\ref{trees}.
These trees are tiny since they are learned   from just the few dozen most
informative examples selected by EZR's active learner.
We use this regression tree format since a recent study by
 Rayegan et al.~\cite{Amiraliminimaldata} showed that these trees can rank features   (via mean decrease in impurity) just as well as other explanation
 algorithms such as  SHAP and ReliefF. Further, when it ranks features,
 EZR needs only  150 labels while SHAP and ReliefF need
 hundreds to thousands of rows.
 
Figure~\ref{htree} shows a software process control task where some manager must 
build a team with (e.g.):
\begin{itemize}
\item
The appropriate programmer capability ({\em pcap});
\item
In order to deliver a system where upper management
is enforcing some   {\em time} crunch on  completion;
\item
While working
in an environment with some level of process maturity ({\em  pmat})
\end{itemize}
Figure~\ref{htree} was built by EZR studying  a data set with 10,000 rows
generated by running what-if queries
on Boehm's COCOMO software models. These models have been extensively
researched and applied, particularly in the United States by the
Department of Defense~\cite{Boehm,boehm2000cost,Boehm:2006,cocomo}.

The left-hand-side of that figure
reports the mean \textsf{disty} of
the rows that fall into different  parts of the tree:
\begin{itemize}
\item
The mean \textsf{disty} of the
untreated rows is 0.52 (see \rc{1}). 
\item The most useful branch of the tree
collects rows whose \textsf{disty} is 0.45
(see \rc{2}).
\end{itemize}
This gap of $0.52-0.45 = 0.07$ looks too
small to be useful. However:
\begin{itemize}
\item
If we look closer at the raw $y$ values,
we see that   even small
changes in \textsf{disty}
can have great business significance.
\item
The
0.52 rows have median DEFECTS=4698, EFFORT=771, KLOC=113.
\item The
0.45 rows have median DEFECTS=4381, EFFORT=379, KLOC=162.
\item That is to say,
the constraints  in the branch down to 
\rc{2}  delivers $(162-113)/113 = 43\%$ more code with
$379/771 = 49\%$ of the effort, with defects roughly unchanged.
 \end{itemize}
\begin{figure}[!t]
{\fontsize{6pt}{7.2pt}\selectfont
\begin{alltt}
disty                                  DEFECTS-   EFFORT-   KLOC- 
-------------------------------------------------------------------
.52                                |    4698  |    771  |  113 \rc{1}
.50    time != xh                  |          |         |
.49    |  pmat != l                |          |         |
.48    |  |  ltex != n             |          |         |
.45    |  |  |  pcap == vh         |    4381  |    379  |  162 \rc{2}
.48    |  |  |  pcap != vh         |          |         |
.48    |  |  |  |  pvol != h       |          |         |
.48    |  |  |  |  |  stor == n    |          |         |
.47    |  |  |  |  |  |  pcap == h |          |         |
.48    |  |  |  |  |  |  pcap != h |          |         |
.48    |  |  |  |  |  stor != n    |          |         |
.50    |  |  |  |  pvol == h       |          |         |
.52    |  |  ltex == n             |          |         |
.51    |  |  |  rely != h          |          |         |
.54    |  |  |  rely == h          |          |         |
.55    |  pmat == l                |          |         |
.48    |  |  rely == n             |          |         |
.60    |  |  rely != n             |          |         |
.54    |  |  |  plex == h          |          |         |
.66    |  |  |  plex != h          |          |         |
\end{alltt}
}
\caption{EZR's regression trees (built using
code from \S\ref{trees}) predict for  \textsf{disty}.
As a method of focusing on essential details, the  $y$ values
are hidden except on two rows.}\label{htree}
\end{figure}
% To evaluate the effectiveness of MDI's explanatory rankings, Rayegan et al.
% ran optimization studies on MOOT data where downstream regressors
% (LR, RF, SVR, ANN, LGBM) were trained using either
% (a)~all attributes, or
% (b)~the top half-dozen attributes as ranked by (b1)~MDI
% or (b2)~two other prominent explanation algorithms:
% SHAP~\cite{lundberg2017unified} and ReliefF~\cite{robnik1997adaptation}.
%  Across 60 datasets, the MDI rankings learned from tiny regression
%  trees were competitive with  SHAP and ReliefF on downstream optimization.
%  But the key result was this: SHAP and ReliefF had to look at all the
%  data while MDI was calculated from a tiny regression tree that only looked
%  at 150 labeled rows.  
% That is, a regression tree grown from a few dozen actively-selected
% examples ranks features about as well as methods that consume the entire
% labeled dataset.
Our tiny regression trees are built from rows selected by an active 
learner. That learner---\textsf{acquire}---relies on two helpers: 
\textsf{warm\_start} (to seed the pool) and a \textsf{score} function 
(to pick what to label next). We describe each, then \textsf{acquire} 
itself.

\textsf{warm\_start} labels \textsf{the.start} random rows 
(default 4), sorts them by \textsf{disty}, and splits them into 
\textit{best} and \textit{rest}. The leftover pool is returned for 
later scoring.
\begin{minted}{python}
the.start = 4 # config

def warm_start(data, rows, label): |\label{warmstart}|
  lab = clone(data, rows[:the.start])
  lab.rows.sort(key=lambda row: disty(lab, label(data, row)))
  n = int(sqrt(len(lab.rows)))
  return (lab, 
          clone(data, lab.rows[:n]),
          clone(data, lab.rows[n:]),
          rows[the.start:])
\end{minted}
EZR allows  customization of the 
\textsf{score} function. For example,
here is a function that ranks unlabeled rows by the difference between their 
distance to the \textit{best} and \textit{rest} centroids. 
\begin{minted}{python}
def acquireWithCentroid(data, best, rest, row): |\label{acqqqq}|
  return (distx(data, row, mids(best)) -
          distx(data, row, mids(rest)))
\end{minted}
In practice, this code is not slow. Recall that 
\cls{Data.mids()} caches centroids until new data invalidates them 
(line~\ref{zap}); hence the centroid
can be  built once  then reused across all unlabeled rows.

(Aside: to stress the main 
message of this paper, observe how the active learning optimizer reuses the 
clustering architecture described above).


\textbf{acquire} stitches everything together. After \textsf{warm\_start}, 
it loops up to \textsf{the.budget} times.
In each loop it uses \textsf{score} to guess  much how unlabeled rows
belong to {\em best} or {\em rest}. These guesses
are used to (a) sort the   unlabeled rows; (b) select the most promising
unlabeled example; (c)~which is then labelled
before (d) it is passed to   \textsf{rebalance}  in order to
u\textit{best} and 
\textit{rest}.
\begin{minted}{python}
the.few = 128   # config
the.budget = 50 # config

def acquire(data, score=acquireWithCentroid, |\label{acquiring}|
                  label=lambda _,row:row): 
  rows = data.rows[:]
  shuffle(rows)
  lab, best, rest, unlab = warm_start(
    data, rows[:the.few], label) |\label{fewing}|
  for _ in range(the.budget):
    if not unlab: break
    pick, *unlab = sorted(unlab,
                     key=lambda row:score(lab,best,rest,row))
    add(lab, add(best, label(data, pick))) |\label{labelling}|
    rebalance(best, rest, lab)
  lab.rows.sort(
    key=lambda row: disty(lab, row))
  return lab |\label{acqured}|
\end{minted}
Three details:
\begin{itemize}
\item
On  line~\ref{fewing}, 
the \textsf{few} caps the number of unlabeled rows searched
 per iteration---larger values did not improve performance.
\item
On line~\ref{labelling},
the 
\textsf{label} function is 
a hook for an oracle  that assigns $y$ values,
For MOOT data, $y$ labels are already present on the row data, so it just returns
the   row. But for other domains, \textsf{label} could (e.g.) call some
simulator to collect new data.
\item
Other active learners have adaptive sampling policies where, early in the run,
they focus on   most uncertain
examples (e.g. where the likelihood of {\em best}) and {\em rest}) is nearly equal). Later in the run, that same learner  might favor  certainty sampling (i.e. find unlabeled examples that are most
likely to be {\em best}). After two years of exploring that kind of adaption, and after extensive
experimentation with the MOOT data, we found   nothing   better than
a
greedy search guided by line~\ref{acqqqq}.
\end{itemize}


 
 \subsection{Testing Active Learning}
The \textsf{acquire} function described above returns a set of labeled
rows. How to test if those conclusions are useful? 

In any realistic 
deployment, a human inspects a short list of recommendations before 
committing to one.
We can evaluate our active learners the same way. This is the standard 
``at-$k$'' practice from information retrieval (precision@$k$, 
hit@$k$, top-$k$ accuracy);
i.e., rather than asking ``did the model rank 
the single best item first?'', we ask ``is a good item somewhere in 
the model's top $k$ picks?''.

In EZR,  $k=$\textsf{the.check} is set to, say, 5.
This \textsf{check} parameter  encodes how much we trust the model: 
\textsf{check}=1 means we trust it enough to just apply its
top recommendation, without checking anything else.
Larger 
values of \textsf{check} are appropriate if we are not 100\% certain about a model.

This  \textsf{check}  is applied in two scenarios:
\begin{itemize}
\item
The \textsf{w1} scenario runs and checks the {\em training data}.
 Given that  we are only allowed to label $M\ll N$ of the rows,
 \textsf{w1} asks,   how good are those $M$ labels at finding the best of the $N$ training
examples? 
\item
\textsf{w2} runs its checks on the {\em testing data}.
Here, we import a model from some past training session and use it
to check some test data. 
\end{itemize}
There is an important difference between these two scenarios.
\begin{itemize}
\item
 Anyone applying the
\textsf{w1} check must do all the active learning, and then all the checking. 
Hence   the \textsf{w1} check requires labeling \textsf{budget}  +  \textsf{check} 
\item
But the \textsf{w2} is far cheaper. Here, some   team has built a model
using some \textsf{budget} which a second team
is  now reusing in a novel situation.
Hence     \textsf{w2} check   just   requires \textsf{check}s.
\end{itemize}
 

\subsubsection{Calculating \textsf{w1} and \textsf{w2}}
To implement those checks, EZR splits each data set in half.
The first half (capped at \textsf{the.few} = 128 rows) goes to
the active learner as the training pool. The second half is
kept separate, and never touched during acquisition.
\begin{minted}{python}
def ready(file):
  d = (file if Data == type(file)
       else Data(csv(file)))
  random.shuffle(d.rows)
  half = len(d.rows) // 2
  return (d, 
          clone(d, d.rows[:half][:the.few]),
          d.rows[half:])
\end{minted}
The controller for this process is \textsf{test\_acquire}. This runs 
20 times, collecting the best row found by the active learner. Each 
repeat (a)~shuffles the data; (b)~gives the first half to the active 
learner; (c)~lets it spend its budget; (d)~grows a tree from the 
labeled rows; and (e)~scores the result using \textsf{w1,w2}  on lines~\ref{w1} 
and~\ref{w2}:
\begin{minted}{python}
the.check  = 5    # config
the.budget = 50   # config

def test_acquire(file=egopt1):
  d0 = Data(csv(file))
  w1, w2 = Num(), Num()
  win = wins(d0) |\label{wins}|
  for _ in range(20):
    d, d_train, test_rows = ready(d0)
    lab = acquire(d_train)
    t   = treeGrow(d_train, lab.rows)
    guess = sorted(test_rows,
              key=lambda r: mid(treeLeaf(t, r).ynum))
    add(w1, win(min( |\label{w1}|
                  lab.rows[:the.check], |\label{check1}|
                 key=lambda r: disty(d_train, r))))
    add(w2, win(min( |\label{w2}|
                 guess[:the.check], |\label{check2}|
                 key=lambda r: disty(d_train, r))))
  return w1,w2
\end{minted}
Since we want to compare results across
many data sets,
our \textsf{win} function used
Equation~\ref{eq:wins}. Further,
to avoid spurious conclusions based on trivially
small distances, we
clamp small improvements to 0. ``Small'' here is Sawilowsky's
threshold of $0.35\sigma$~\cite{sawilowsky2009new}, with $\sigma$
estimated robustly via  the standard method of   (90th $-$ 10th)/2.56 percentile range
of \textsf{disty}.
\begin{minted}{python}
def wins(data): |\label{wins}| 
  ys = sorted(disty(data, row1) for row1 in data.rows)
  ten = len(ys) // 10
  lo  = ys[0]
  med = ys[5*ten]
  sd  = (ys[9*ten] - ys[ten]) / 2.56
  def f(row2):
    x = disty(data, row2)
    if x < lo + 0.35*sd: x = lo
    return max(-100,
          int(100 * (1 - (x - lo)/(med - lo + 1e-32))))
  return f
\end{minted}

 
\begin{figure*}
\noindent
\begin{tabular}{@{}>{\raggedleft\arraybackslash}m{0.35\textwidth}@{\hspace{1em}}m{4in}@{}}
\textsf{w1} results (test on training set):\newline
(all values computes using line~\ref{wins})\newline
  & \includegraphics[width=4in]{train.pdf} \\[1ex]
\textsf{w2} results (test on a separate hold-out set):\newline
(all values computes using line~\ref{wins})\newline
  & \includegraphics[width=4in]{test.pdf}
\end{tabular}
\caption{\textsf{w1} and \textsf{w2} results in 100,000 trials.
For each trial:
(a)~pick a MOOT data set at random;
(b)~pick \textsf{budget} and \textsf{check} at random from
$1 \le B \le 150$ and $1\le C \le 10$;
(c)~apply active learning report the wins
calculated from line~\ref{wins}.
 Contour lines mark \textsf{wins} of 60, 75,
85, 95. Each cell aggregates the median over many trials
at that (\textsf{budget}, \textsf{check}) pair, smoothed by a
small Gaussian blur from neighboring cells
(where cells $c_1,c_2,c_3$ at distance
1, 2, 3 contribute a weight $(71c_1 + 25c_2 + 4c_3)/100$).
 Contours flatten
toward the upper right, indicating diminishing returns from
larger budgets.}
\label{check}
\end{figure*}


\subsubsection{Parameter Effects}
Using the \textsf{w1} and \textsf{w2} scores,
we can comment on the validity of   EZR's default active
learning parameters of  \[\textsf{budget},\textsf{check}) = (50,5)\]
In this section, we test if better settings (e.g. a larger training budget)
would lead to fundamentally better results.

 Figure~\ref{check} presents the results of a study where   MOOT data sets were selected
at random, then explored with \textsf{budget},\textsf{check} between 
  $[1,150]$ and  $[1,10]$. This process was repeated    100{,}000 times.

The top panel of Figure~\ref{check} reports \textsf{w1}: i.e. 
it is  scored on rows drawn
from its own training pool. 
The bottom panel of Figure~\ref{check} reports the  \textsf{w2}
results where we score on
  held-out rows     not seen during acquisition.
Recall from the above discussion that \textsf{w2} occurs when
a second team is reusing a model developed by another team
(and in  this \textsf{w2} scenario, this second team
only has to label \textsf{check} number of rows).

There are several noteworthy effects in     Figure~\ref{check}:
\begin{itemize}
\item The \textsf{w1} scores
are less than the \textsf{w2} scores. This is to be expected:
testing on training data tends to overestimate the performance on as-yet-unseen future testing data. 
\item At EZR's default values of (\textsf{budget},\textsf{check}) = (50,5),
we see our active learners usually achieve above an  85\% \textsf{win}.
That is to say, reusing someone else's model and only labeling five new examples
is enough to get most of the way to the reference optimum. 
\item
If we increase either of  (\textsf{budget},\textsf{check}) = (50,5), then we move to the top-right
of the bottom panel. In that region, we can see improvement of    the performance score on hold-out data.

\item
That said, the improvements above (50,5) are modest.   As we move towards the top right of the bottom panel, the performance contours start flattening
out as  performance grows to 100 (which is the maximum score possible on these data sets).
\end{itemize}
From this we conclude that for safety-critical and mission-critical applications,
it can be useful to move beyond (\textsf{budget},\textsf{check}) = (50,5): specifically,
{\em more checks} seem more useful than {\em more training budget}. And checking beyond (say) 7 items
seems not particularly useful. 
But for many engineering applications, EZR's current default settings would  suffice. 

\begin{figure*} 
\begin{center}
\includegraphics[width=\linewidth]{used.pdf}
\end{center}
\caption{Across 120+  MOOT tasks, the
\color{darkyellow}{\bf yellow line}\color{black}~shows the
number of $x$ attributes available (rising to over a thousand).
The \color{red}{\bf red line}\color{black}~shows the number of
unique attributes \IT{}'s trees actually used (under ten,
throughout). The \color{blue}{\bf blue line}\color{black}~shows
the \textsf{w2} (calculated using \textsf{wins}   from line~\ref{wins}) of those trees,
hovering near 100 across the entire range. Note that the  blue line does
{\bf not} slope downward from left to right.
That is, it is not true that performance drops if  EZR ignores more and more of the attributes.}
\label{used}
\end{figure*}

 ~\\~\\
\subsubsection{Model Size}

 
Returning to Figure~\ref{htree}, that tree 
had an interesting feature that was not previously discussed. 
That tree splits its data on 8 different  attributes.
This is noteworthy since that data  set has 24 unique attributes;
i.e. EZR  built that model by ignoring most of the attributes.


This is not a one-off result: in the usual case,
EZR ignores most of the attributes.
Figure~\ref{used} shows how many attributes were used
in the tiny regression trees built after active
learning on   120+ MOOT tasks. 
 For example, in the {\em auto93} data set,
the goal is to maximize miles-per-gallon and acceleration without
increasing the weight of the car.

For these optimization tasks, Figure~\ref{used} shows three lines:

\begin{itemize}
\item The \color{darkyellow}{\bf yellow line}\color{black}~ is
the number of $x$ attributes available in each task. It rises
from a few dozen to over a thousand as we move right across
MOOT.
\item The \color{red}{\bf red line}\color{black}~ is the number
of unique attributes that actually appear in any of the trees
\IT{} grew on each task. It stays under ten throughout.
\item The \color{blue}{\bf blue line}\color{black}~ is the
\textsf{w2} 
\textsf{wins} score (from line~\ref{wins}) of those trees. 
This line  hovers near 100 across the entire range.
\item
This  \color{blue}{\bf blue line}\color{black}~line does not slope downward as we move from left to right. 
That is, 
for the  120+ MOOT tasks, fewer than ten variables suffice to
build effective predictive models, even when the data set
contains hundreds or thousands of attributes.
\end{itemize}
This {\em feature selection} result
has a long history in AI. Pearson described principal
components in 1902~\cite{pca}. Chang showed in 1974 that a
handful of exemplar rows can replace a full nearest-neighbor
training set without loss of classification
accuracy~\cite{chang74}. Johnson and Lindenstrauss proved in
1984 that high-dimensional data can be projected to far fewer
dimensions while preserving pairwise
distances~\cite{johnson1984extensions}. Kohavi and John reported
in 1997 that 80\% of features can usually be
discarded~\cite{kohavi97}. Zhu's 2005 semi-supervised learning
survey describes high-dimensional data as typically lying on a
much lower-dimensional manifold~\cite{zhu2005semi}. Settles
argued in 2009 that an active learner querying only the most
informative rows can match a learner trained on the full
set~\cite{settles2009active}. Our own work on SE ``keys'' across
2003--2021 found that a handful of parameters governs the
behavior of many software
models~\cite{me03a,me07a,me21a,menzies2008implications}.

What is new in \IT{} is that we extends this feature selection effect
to
multi-objective optimization. Most of the prior work cited above
was about classification or regression. \IT{}, on the other hand, grows trees that
predict \textsf{disty}, which means the trees are doing
multi-objective ranking. The same effect (under ten variables
suffice) holds in that setting too.

\subsubsection{Comparison with State-of-the-Art}\label{sota}

In other publications, we have compared the \textsf{w2} scores generated
 with state-of-the-art explanation and optimization algorithms.


Rayegan et al.~\cite{Amiraliminimaldata} compared \IT{}'s trees against
standard AI explanation algorithms
(SHAP, LIME, and Anchors) on standard XAI benchmarks. The metrics collected
were {\em actionability} (does the explanation tell you what to
change?) and {\em faithfulness} (does following the explanation
actually move the prediction?). \IT{}'s trees matched or
exceeded the named baselines on both metrics.
This is a useful property in safety-relevant settings. A
practitioner reading a SHAP report cannot easily tell whether
the surrogate has departed from the underlying model. Reading a
\IT{} tree, the practitioner is reading the actual decision
boundary. If the tree says \texttt{pcap == vh} buys 49\% effort
reduction, then \texttt{pcap == vh} buys 49\% effort reduction.

Ganguly et al.~\cite{ganguly2026lowgodatalightse}    compare \IT{}'s
optimizations against state-of-the-art optimizers including SMAC3~\cite{hutter2011sequential}
and others.
As mentioned above, they found that  EZR ran 500 times faster than SMAC3. 
They also found that EZR delivered competitive results using a fraction
of the data needed by other methods.

So, in summary,
on the MOOT data, a  fewer labels,
fewer features, and lighter model updates each match or beat
more complete alternatives. And the software that achieves all
this is a very small extension to the EZR substrate used
for classification, regression and optimization.


\section{Text Mining}\label{sec:textmine}

Some of the above results have been published before (e.g. all the results
in \S\ref{sota}). 
Hence, to cap off this paper,
 we offer a new text mining result. Specifically, we successfully
 apply all the above machinery
 to text mining and relevancy filtering of documents.
 
 In the current literature,  such text mining is often
 explored using a RAG-like approach (retrieval-augmented
 generation, described in the next section). 
This requires LLM-grade embedding spaces, billion-parameter encoders,
and similarity queries against a vector database. 

Older
approaches explored text mining and relevancy filtering using
support vector machines. ~\cite{yu2018total,yu2018finding}. This section shows
that 30 lines of complementary Bayes~\cite{rennie2003tackling}
on the substrate from \S\ref{sec:peek}. These results
are so good  that there is little point
in exploring other methods (e.g. LLMs). Hence we have to ask:
just how simple is text mining? Perhaps it it time
to conduct more experiments
comparing LLMs to simpler alternatives.



 \subsection{Retrieval-Augmented Generation} 
 RAG~\cite{lewis2020rag} was a large language model
 technique that was used extensively when LLMs
 had smaller context
 windows. In order to avoid exhausting the space
 inside a context window,
   queries to an LLM were 
augmented (on a query-by-query basis)
with a small number of relevant chunks
from (e.g.)  documents stored on disk. 

A RAG system has two halves: a post-query dialog generator and a
 search method for relevant documents. While EZR (as yet)
has no dialog generation method, we show below
that finding relevant documents is surprisingly easy using
a Naive Bayes trick built
around EZR.

\subsection{Yu et al. and FASTREAD}

In this section, we compare our results to prior work 
by
Yu et al.~\cite{yu2018total}.
That work
explored methods for 
finding relevant text documents.
They explored  
{``primary study selection'' for
systematic literature reviews (SLR)
in software engineering. 
That is, the goal here is 
to read
the research literature
to find the $M\ll N$
papers that are most relevant
to some research query.
 Reading
every candidate takes weeks or months~\cite{yu2017,yu2018}. 
To reduce the number of papers that must be read, Yu's tool 
(called FASTREAD) used a linear SVM with
some sampling tricks:
\begin{itemize}
\item 
When initializing his SVM classifiers for {\em positive examples},
they began with a small number of ``warm start'' papers offered by domain experts. 

\item To
find a set of initial {\em negative examples}, Yu used {\em presumptive sampling}:
  an equal number of randomly-selected papers, declared negative on
  the assumption that, when the target class is rare, a random draw
  is overwhelmingly unlikely to belong to it.
  \item After initialization, Yu et al.\ used their own active learning
  methods to incrementally decide what to sample next.
\end{itemize}
Their approach yielded a 95\% recall on Yu's four corpora after labeling
roughly 300--800 papers out of 1,700--8,900. 
Yu et al.\ report that that was a
20--50\% improvement over the prior state of the
art.

For the rest of this paper,  we  ask if EZR can perform the same task using, say,
Naive Bayes and the {\em best} and {\em rest} models
used in our active learning.  
In this study ``best'' means most relevant to the human conducting an SLR.

\subsection{Rennie et al. and Complementary Bayes}
Standard classifiers like Naive Bayes have problems with text mining. 
Rennie et al.~\cite{rennie2003tackling} diagnose the problem as one
of ``bursty'' low frequency counts inherent in the sparse data used to train text miners.
Relevancy filtering in text mining is a particularly challenging low frequency problem since, within 10,000s of documents, there may only be (say) 50 relevant documents;
i.e., we are asking a classifier to predict for a class with a one-in-200 frequency (or rarer).  


Rennie et al.'s approach to this problem was 
to invert the normal classification algorithm:
\begin{itemize}
\item
When training their {\em complementary Bayes classifier}, a new row updates the \cls{Data} distributions
for all the classes to which it does {\em not} belong. 
\item
Then, when classifying, 
instead of
reporting what class is most likely,    {\em complementary Bayes}
reports the class that a row is least likely to {\em not} belong.
\end{itemize}
We inserted this complementary Bayes classifier into our active learner as follows.
Recall that 
the \textsf{acquire} function of line~\ref{acquiring} accepts a \textsf{score} function
that guesses which unlabeled row is more likely to be  {\em best}. The following relevancy filter  hooks into the likelihood calculation of complementary Bayes.
\begin{minted}{python}
def acquireWithBayes(data, best, rest, row): 
  n = len(best.rows) + len(rest.rows)
  return likes(rest, row, n, 2) - likes(best, row, n, 2)
\end{minted}
We tested if this active learner approach can solve
the same problems   explored by Yu et al.\ (see Table~\ref{slr}).
Yu's test cases  come from 
four real-world published SE SLRs:
\begin{itemize}
\item
Hall et al.'s work~\cite{hall2012} on fault
prediction;
\item
Wahono~\cite{wahono2015systematic}'s study  on defect prediction;
\item
Radjenovi\'c et al.'s study~\cite{radjenovic2013software} on software fault
metrics;
\item
Kitchenham et al.'s study~\cite{kitchenham2013systematic} on a
tertiary review of SLR practice.
\end{itemize}
In each case, these humans (all expert
in the field of SE) read thousands of papers to select a few dozen
relevant to their research. 
The data sets of Table~\ref{slr} were extracted from the careful 
logs
they kept of their reading process. 

\begin{figure*}[!t]
\centering
\includegraphics[width=\linewidth]{kit.pdf}
\caption{Recall (top row) and false alarm (bottom row) as a
function of labels acquired, for non-normalized TWCNB (blue)
and normalized TWCNB (red) on Yu's four SLR corpora. Warm-start
= 24 for all runs. Each curve is the median over 20 repeats,
with shaded IQR. Vertical bars on the recall panels mark the
labeling budget at which Yu et al.~\cite{yu2018finding} reached
the same recall with FASTREAD. Numerical labels at the left
edge of each recall panel mark the recall achieved
immediately after the warm start.}
\label{slrfig}
\end{figure*}
 

\begin{table}[!t]
\caption{Yu's four SLR corpora~\cite{yu2018total}. Each is a
search-engine query result where a small minority of candidates
are relevant. The Kitchenham corpus has two relevance levels:
132 papers passed title-and-abstract review, of which 45
survived content review.}
\label{slr}
\begin{center}
\footnotesize
\begin{tabular}{lrr}
\hline
Dataset & \#Candidate & \#Relevant \\
\hline
Wahono       & 7002 & 62 \\
Hall         & 8911 & 104 \\
Radjenovi\'c & 6000 & 48 \\
Kitchenham   & 1704 & 132 \\
\hline
\end{tabular}
\end{center}
\end{table}

\subsection{Results}
To assess these results, we tracked
 two metrics:
 \begin{itemize}
 \item

\textit{Recall} is $\mathit{tp}/(\mathit{tp} + \mathit{fn})$:
the fraction of truly relevant papers found. 
{\em Higher} recalls are {\em better} since it means we are finding
more relevant papers.
  Yu's study achieved a recall of  95\%.
\item
\textit{False alarm} (sometimes called pf, the false positive
rate) is $\mathit{fp}/(\mathit{fp} + \mathit{tn})$: the fraction
of irrelevant papers incorrectly returned as relevant. 
{\em Lower} false alarms are {\em better} since, otherwise,
our human experts are wasting time reading irrelevancies.
Yu's study did not report false alarm.
\end{itemize}
 


\definecolor{slateblue}{RGB}{106,142,162}

\definecolor{ruby}{RGB}{177,98,143}




 Figure~\ref{slrfig} shows results of 20 repeated trials where,
each time, the order of the training rows is shuffled
(in this case, each ``row'' was the title and abstract of a different paper).
Each run was initialized with a warm start of $N=24$ randomly drawn papers. 
This number was selected after asking academics in our field ``for any particular area, how many dozen papers do you usually
know?''

After that,
complementary Bayes selected one paper at a time (the unlabeled paper
it ranked most strongly as ``yes''), and recall and false alarm were
updated after each new label. The total labeling budget was 50. 


 

In  Figure~\ref{slrfig}, 
solid lines show the 50\% percentile, and the shaded
regions show the 25th to 75th range:
\begin{itemize}
\item
The four columns show results from our four 
corpora.
\item
The top row
of plots shows recall.
\item
The second row shows  false alarm rates
\item The 
 \color{ruby}{\bf red lines}\color{black}~ shows results using
 Rennie et al.'s recommendation of normalizing a likelihood
 by dividing it by the sum of all likelihoods.
 \item
 The  \color{slateblue}{\bf blue lines}\color{black}~
 shows results where normalization was disabled.
 \end{itemize}
Note that the false alarm rate is far lower than
 when normalization is disabled; i.e., that part of
 the Rennie et al.\ result should not be applied
 uncritically to new data.
Clearly,  measuring {\em both} recall and false
alarms is informative for relevancy filtering.
 Yu et al.\ did not measure false alarms, and we would now
 call that a methodological error.
  
The most striking result here is the early plateaus.
High recalls and low false alarms are possible in all our 
datasets, and these can be achieved after 100 examples or less. By comparison, Yu et al.\ required 300 to 600   labels to reach the
same recall. On Kitchenham, the recall plateaus near 82\%, below
the 95\% target, but complementary Bayes still recovers 108 of the 132 relevant
papers within 40 labels. For an SLR conductor whose goal is to
find relevant papers quickly, that is a useful result.

The contrast with FASTREAD is sharper than the absolute numbers
suggest. Yu's recall curves begin at very low values and climb
only about 7\% across several hundred labels. Ours begin near
the ceiling and remain there. Same data, same task, very
different rates of convergence.

In summary, compared to Yu's methods, EZR+complementary Bayes:
\begin{itemize}
\item Needs fewer labels (100 vs.\ 300+);
\item Uses a simpler classifier (Naive Bayes rather than a
  linear SVM);
\item Uses a smaller codebase (based on EZR, complementary
Bayes is just three dozen lines of code);
\item Reports both recall and false alarm, which exposed the
  finding that one of Rennie's recommended preprocessing
  steps is counterproductive in this domain.
\end{itemize}

That last item carries the methodological lesson of this
section. Yu's data was fine, but the choice of metrics was
incomplete, as revealed by an investigation with  simpler tools
  \IT{}).
We suspect that
similar gaps exist elsewhere in the empirical SE literature,
masked by tools too slow to encourage re-checks. Compact
toolkits like EZR make those re-checks cheap, and we suggest
the community should run them.

A broader point bears on the main argument of this paper.
Relevancy filtering in RAG is increasingly framed as a task for large
language models, using embedding spaces of billions of
parameters and similarity queries against a vector database.
Our experiment achieved comparable or better results with 30
lines of complementary Bayes riding on the same
\cls{Num}/\cls{Sym}/\cls{Data} substrate used throughout
EZR. We wonder how many other methods from the LLM community could be greatly simplified by simpler methods.

% In summary (a)~relevance filtering may not always need
% billion-parameter embeddings nor engineered linear SVMs since
% (b)~thirty
% lines of complementary Bayes on a shared substrate can work surprisingly well.

% --- end of section ---

\section{Summary and Discussion}\label{sec:disc}

This paper set out to defend the claims of
Table~\ref{truisms}: that new problems need not always be
tackled with more complex algorithms and more data. We
offered 400 lines of Python that:
\begin{itemize}
\item Built a simple substrate (\cls{Data}, \cls{Num},
\cls{Sym}, and \textsf{add}) in \S\ref{sec:peek};
\item Used that substrate to implement, in just a few lines,
Naïve Bayes, $k$-means, and $k$-means++ (\S\ref{sec:famous});
\item Then classification and regression trees
(\S\ref{trees});
\item Then two optimizers, simulated annealing and local
search (\S\ref{sec:opt});
\item Then an active learner (\S\ref{sec:active});
\item And a text-mining relevance filter
(\S\ref{sec:textmine}).
\end{itemize}

Along the way we found that:
\begin{itemize}
\item In its 1983 default configuration, simulated annealing
can beat 1990s local-search variants --- despite a decade of
papers claiming the newer variants had surpassed it
(\S\ref{sec:opt});
\item Our active learner reached the reference optimum with
fewer than 100 labels and fewer than 10 features (even when
thousands of features were available), and ran 500$\times$
faster than state-of-the-art algorithms like SMAC3
(\S\ref{sec:active});
\item Text-mining tasks that are now routinely handed to
massive language models can be implemented by 30 lines of
complementary Bayes on the substrate above, and the same
experiment exposed a measurement gap that the prior
heavyweight evaluation had missed (\S\ref{sec:textmine}).
\end{itemize}
Just to be clear, we have not shown 
that  all SE problems can be simplified using \IT{}.
This study has a clear sample bias. The six counter-findings
of Table~\ref{truisms} were extracted from experience with
the MOOT data, which holds examples of a specific family of
SE problems:
\begin{itemize}
\item Structured tabular data with explicit $x$ and $y$
columns, where $y$ is expensive to obtain.
\item That family is very broad (see Table~\ref{mootdata})
and includes configuration tuning, performance prediction,
defect prediction, test selection, cost estimation, and (as
\S\ref{sec:textmine} shows) some text-mining tasks.
\item But broad is not universal. Some problems are
fundamentally complex. The certification of safety-critical
software is not (and should not be) a simple process.
Generation tasks (translation, summarization,
natural-language interfaces, code synthesis) genuinely need
the data scaling that LLMs provide.
\end{itemize}
 Another bias in this study is how we measured success. Our
own experiments reported medians over 20 or more repeats
using the \textsf{wins} score (Equation~\ref{eq:wins}), which
clamps trivially small differences ($<0.35\sigma$) to zero
per Sawilowsky~\cite{sawilowsky2009new}. State-of-the-art
comparisons against SHAP, LIME, Anchors, SMAC3, TPE, and
HyperOpt come from independent replications by Rayegan et
al.~\cite{Amiraliminimaldata} and Ganguly et
al.~\cite{ganguly2026lowgodatalightse}, using their own
methodologies. By these criteria, every comparison we ran
--- on 120+ MOOT tasks, on Yu et al.'s four SLR
corpora~\cite{yu2018total}, and in the independent
replications --- showed the tiny EZR-based approach running
faster and performing as well as or better than the heavier
state-of-the-art tool. A different metric could change which
method appears to win, and we cannot rule that out.

The reverse statement is also worth making. The SE empirical
literature has many examples where the published advantage
of a heavyweight method shrank or disappeared once re-tested
against a carefully built simple baseline~\cite{Fu17,
fu2016tuning,agrawal2019dodge,Nair17a,
tantithamthavorn2016automated,grinsztajn2022why,Hou24}.
\IT{} is one more entry in that pattern.

In conclusion, just because some tasks are hard does not
mean all tasks are hard. The challenge we put to the
community is this: have we really checked what is genuinely
complex, and what is, on careful reading, actually very
simple?

\bibliography{bib3,main}
\end{document}
